%% toc: P2 19 19.1,19.2,19.3
\section{Proofs for Herbrand's theorem, the clause form and Statman's direct computations}
\label{sec:annex-p02-w02-p02-028}

Throughout the first part, the standing assumptions are those of
Convention~\ref{def:p02-w02-p02-028-standing}.  ``Structure'' means a
structure with nonempty domain.  We use one elementary fact about reducts
twice: \emph{if \(\Lang'\subseteq\Lang\), \(\mathcal N\) is an
\(\Lang\)-structure and \(\psi\) is an \(\Lang'\)-formula, then
\(\mathcal N,s\entails\psi\) iff \(\mathcal N\restriction\Lang',s\entails\psi\)}.
This holds by induction on terms and then on formulas, since every clause of
the definitions of term value and satisfaction consults only the domain, the
assignment and the interpretations of symbols occurring in the expression
evaluated.  Conversely, every \(\Lang'\)-structure \(\mathcal N'\) is the
reduct of an \(\Lang\)-structure: fix one element \(d\) of its domain and
interpret each constant of \(\Lang\setminus\Lang'\) by \(d\), each function
symbol by the constant function with value \(d\), and each relation symbol by
the empty relation.

\begin{annexitem}{annex:p02-w02-p02-028-ground-evaluation}{thm:p02-w02-p02-028-ground-evaluation}{Ground evaluation}
\begin{proof}
(1) Induction on \(\psi\).  For a ground atom \(A\),
\(\mathcal N\entails A\) iff \(v_{\mathcal N}(A)=1\) by the definition of
\(v_{\mathcal N}\).  \(\bottom\) is false in \(\mathcal N\) and has value \(0\).
For \(\neg\psi_1\), \(\psi_1\land\psi_2\), \(\psi_1\lor\psi_2\) and
\(\psi_1\to\psi_2\), the satisfaction clause for a sentence and the
truth-table clause give the same truth value from the same values of the
components, which agree by the induction hypothesis.  Assignments play no
role because \(\psi\) has no variables (Lemma~\ref{lem:fol-coincidence}).

(2) Put \(s'=s[x_1\mapsto t_1^{\mathcal N},\ldots,x_n\mapsto t_n^{\mathcal N}]\)
and write \(r(\bar t)\) for the simultaneous replacement of each \(x_i\) by
\(t_i\) in a term \(r\).  We first show
\(\sem{r(\bar t)}^{\mathcal N}_s=\sem{r}^{\mathcal N}_{s'}\) by induction on
\(r\).  If \(r\) is \(x_i\), the left side is \(t_i^{\mathcal N}\), the value
of a ground term, which is \(s'(x_i)\).  If \(r\) is another variable \(y\),
both sides are \(s(y)=s'(y)\).  If \(r\) is a constant, both sides are its
interpretation.  If \(r=f(r_1,\ldots,r_k)\), both sides are
\(f^{\mathcal N}\) applied to values that agree by the induction hypothesis.
For an atom \(R(r_1,\ldots,r_k)\) the claim
\(\mathcal N,s\entails R(r_1,\ldots,r_k)(\bar t)\) iff
\(\mathcal N,s'\entails R(r_1,\ldots,r_k)\) follows from the term equation,
and the connective cases follow from the induction hypothesis, because the
replacement commutes with each connective.  Since \(\varphi(\bar t)\) is a
sentence, its truth does not depend on \(s\).

(3) For a constant \(d\), \(d^{\mathcal M_B}=d\) by definition.  If
\(t_i^{\mathcal M_B}=t_i\) for each \(i\), then
\(f(t_1,\ldots,t_n)^{\mathcal M_B}=f^{\mathcal M_B}(t_1,\ldots,t_n)=f(t_1,\ldots,t_n)\).
Hence for a ground atom,
\(\mathcal M_B\entails R(t_1,\ldots,t_n)\) iff
\((t_1^{\mathcal M_B},\ldots,t_n^{\mathcal M_B})=(t_1,\ldots,t_n)\) satisfies
\(R^{\mathcal M_B}\), iff \(R(t_1,\ldots,t_n)\in B\); for a nullary \(R\),
iff \(R\in B\).  So \(v_{\mathcal M_B}=v_B\).
\end{proof}
\end{annexitem}

\begin{annexitem}{annex:p02-w02-p02-028-herbrand-model}{thm:p02-w02-p02-028-herbrand-model}{Herbrand model lemma}
\begin{proof}
Since \(\Lang\) is countable, its ground terms are finite strings over a
countable alphabet, so the Herbrand universe \(U\) and the Herbrand base are
countable; \(U\) is nonempty because \(\Lang\) has a constant.  Fix an
injection \(\iota\) of the Herbrand base into \(\{p_n:n\in\mathbb N\}\), and
for a ground quantifier-free \(\psi\) let \(\psi^\iota\in\mathsf{Fm}_{\mathsf{CPL}}\) replace
each ground atom \(A\) by \(\iota(A)\) and keep every connective and
\(\bottom\).  For any propositional valuation \(w\), induction on \(\psi\)
shows that \(\psi^\iota\) has \(w\)-value \(\llbracket\psi\rrbracket_{w\circ\iota}\).

Let \(\Gamma^\iota=\{\gamma^\iota:\gamma\in\Gamma\}\).  A finite subset of
\(\Gamma^\iota\) is \(\Gamma_0^\iota\) for a finite \(\Gamma_0\subseteq\Gamma\).
By hypothesis some atom valuation \(v_0\) gives every member of \(\Gamma_0\)
the value \(1\).  Define \(w_0(\iota(A))=v_0(A)\) and \(w_0(p)=0\) for \(p\)
outside the range of \(\iota\); then \(w_0\circ\iota=v_0\), so \(w_0\)
satisfies \(\Gamma_0^\iota\).  By Theorem~\ref{thm:prop-compactness}(2) some
valuation \(w\) satisfies all of \(\Gamma^\iota\).  Put
\(B=\{A:w(\iota(A))=1\}\), so \(v_B=w\circ\iota\) and every
\(\gamma\in\Gamma\) has \(\llbracket\gamma\rrbracket_{v_B}=1\).  By
Lemma~\ref{thm:p02-w02-p02-028-ground-evaluation}(3) and~(1),
\(\mathcal M_B\entails\gamma\) for every \(\gamma\in\Gamma\).

Now let \(\varphi(x_1,\ldots,x_n)\in\Phi\) and let \(s\) be any assignment in
\(\mathcal M_B\).  Each \(s(x_i)\) is a ground term \(t_i\), and
\(t_i^{\mathcal M_B}=t_i\) by
Lemma~\ref{thm:p02-w02-p02-028-ground-evaluation}(3), so \(s\) agrees on
\(x_1,\ldots,x_n\) with \(s[x_1\mapsto t_1^{\mathcal M_B},\ldots]\).  By
Lemma~\ref{lem:fol-coincidence} and
Lemma~\ref{thm:p02-w02-p02-028-ground-evaluation}(2),
\(\mathcal M_B,s\entails\varphi\) iff \(\mathcal M_B\entails\varphi(\bar t)\),
which holds because \(\varphi(\bar t)\in\Gamma\).  As \(s\) was arbitrary,
\(\mathcal M_B\) satisfies the universal closure of \(\varphi\).
\end{proof}
\end{annexitem}

\begin{annexitem}{annex:p02-w02-p02-028-existential-validity}{thm:p02-w02-p02-028-existential-validity}{Herbrand's theorem for existential sentences}
\begin{proof}
Write \(\sigma=\exists\bar x\,\varphi\).  When \(n=0\) the only expansions
are \(\varphi\lor\cdots\lor\varphi\), and both directions below specialize
without change.

\emph{Right to left.}  Suppose
\(D=\varphi(\bar t_1)\lor\cdots\lor\varphi(\bar t_m)\) is propositionally
valid and let \(\mathcal N\) be any \(\Lang\)-structure.  Then
\(\llbracket D\rrbracket_{v_{\mathcal N}}=1\), so
\(\mathcal N\entails D\) by
Lemma~\ref{thm:p02-w02-p02-028-ground-evaluation}(1), and
\(\mathcal N\entails\varphi(\bar t_i)\) for some \(i\).  By
Lemma~\ref{thm:p02-w02-p02-028-ground-evaluation}(2), for any assignment
\(s\), \(\mathcal N,s[\bar x\mapsto\bar t_i^{\mathcal N}]\entails\varphi\),
and \(n\) applications of the satisfaction clause for \(\exists\) give
\(\mathcal N,s\entails\sigma\).  Hence \(\sigma\) is valid.

\emph{Left to right.}  Suppose \(\sigma\) is valid.  Let \(\Lang_\sigma\)
consist of the constant, function and relation symbols occurring in
\(\sigma\), together with one constant \(c_0\) of \(\Lang\) when \(\sigma\)
contains no constant.  \(\Lang_\sigma\) is finite.

\emph{Step 1: \(\sigma\) is true in every \(\Lang_\sigma\)-structure.}  An
\(\Lang_\sigma\)-structure is the reduct of some \(\Lang\)-structure, which
satisfies \(\sigma\); by the reduct fact, so does the
\(\Lang_\sigma\)-structure.

\emph{Step 2.}  Suppose, for a contradiction, that no expansion
\(\varphi(\bar t_1)\lor\cdots\lor\varphi(\bar t_m)\) with ground
\(\Lang_\sigma\)-terms is propositionally valid.  Let \(\Phi=\{\neg\varphi\}\)
and let \(\Gamma\) be the set of its ground \(\Lang_\sigma\)-instances
\(\neg\varphi(\bar t)\).  For a nonempty finite subset
\(\{\neg\varphi(\bar t_1),\ldots,\neg\varphi(\bar t_m)\}\), the disjunction
\(\varphi(\bar t_1)\lor\cdots\lor\varphi(\bar t_m)\) is not propositionally
valid, so some atom valuation gives every \(\varphi(\bar t_i)\) the value
\(0\) and every \(\neg\varphi(\bar t_i)\) the value \(1\); the empty subset
is trivially satisfiable.  By
Lemma~\ref{thm:p02-w02-p02-028-herbrand-model}, applied to the countable
signature \(\Lang_\sigma\), some Herbrand \(\Lang_\sigma\)-structure
\(\mathcal M_B\) satisfies \(\forall\bar x\,\neg\varphi\), hence does not
satisfy \(\sigma\).  This contradicts Step~1.

Hence some expansion with ground \(\Lang_\sigma\)-terms is propositionally
valid.  Since \(\Lang_\sigma\subseteq\Lang\), its terms are ground
\(\Lang\)-terms and it is a finite Herbrand expansion of \(\sigma\) over
\(\Lang\), whose terms use only the symbols stated.
\end{proof}
\end{annexitem}

\begin{annexitem}{annex:p02-w02-p02-028-duality}{thm:p02-w02-p02-028-herbrand-skolem-duality}{Herbrand form as negated Skolem form}
\begin{proof}
\(\sigma^{\neg}\) has the same prefix variables as \(\sigma\), each bound
once, and quantifier-free matrix \(\neg\theta\), so it is a rectified prenex
sentence.  For every structure, assignment and formula \(\rho\),
\(\mathcal N,s\entails\neg\forall x\,\rho\) iff not every \(d\) in the domain
gives \(\mathcal N,s[x\mapsto d]\entails\rho\), iff some \(d\) gives
\(\mathcal N,s[x\mapsto d]\entails\neg\rho\), iff
\(\mathcal N,s\entails\exists x\,\neg\rho\); dually
\(\neg\exists x\,\rho\) and \(\forall x\,\neg\rho\) are equivalent.  If
\(\rho\) and \(\rho'\) are equivalent, so are \(Qx\,\rho\) and
\(Qx\,\rho'\).  Induction on the length of the prefix now gives
\(\neg Q_1x_1\cdots Q_nx_n\,\theta\) equivalent to
\(\overline Q_1x_1\,\neg Q_2x_2\cdots Q_nx_n\,\theta\) and then to
\(\sigma^{\neg}\).

Skolemizing \(\sigma^{\neg}\) treats exactly the positions \(i\) with
\(Q_i=\forall\), which are existential in \(\sigma^{\neg}\).  The universal
positions of \(\sigma^{\neg}\) before such an \(i\) are the existential
positions of \(\sigma\) before \(i\), whose variables are the \(u_1,\ldots,u_m\)
of Definition~\ref{def:p02-w02-p02-028-herbrand-form}.  Hence, with the
Skolem symbol at \(i\) named \(h_i\), the Skolem term at \(i\) is
\(h_i(u_1,\ldots,u_m)\), the retained quantifiers are
\(\forall u_1\cdots\forall u_p\), and the matrix becomes \(\neg\theta\) with
each \(x_i\) (for \(Q_i=\forall\)) replaced by its term.  Replacement in a
quantifier-free formula commutes with \(\neg\), so this matrix is
\(\neg\theta^{\mathrm H}\).
\end{proof}
\end{annexitem}

\begin{annexitem}{annex:p02-w02-p02-028-hf-validity}{thm:p02-w02-p02-028-herbrand-form-validity}{Validity is preserved by the Herbrand form}
\begin{proof}
\emph{First implication, without choice.}  Suppose \(\sigma\) is valid but
\(\operatorname{Hf}(\sigma)=\exists\bar u\,\theta^{\mathrm H}\) is not.  Then
some \(\Lang^{\mathrm H}\)-structure \(\mathcal N\) satisfies
\(\neg\exists\bar u\,\theta^{\mathrm H}\), which is equivalent to
\(\forall\bar u\,\neg\theta^{\mathrm H}=\operatorname{Sk}(\sigma^{\neg})\) by
the quantifier dualities used in the previous proof and
Lemma~\ref{thm:p02-w02-p02-028-herbrand-skolem-duality}.  By
Theorem~\ref{thm:skolem-equisatisfiable}(2), which uses no choice, the
\(\Lang\)-reduct of \(\mathcal N\) satisfies \(\sigma^{\neg}\), hence
\(\neg\sigma\).  This contradicts validity of \(\sigma\).

\emph{Second implication, with the Axiom of Choice.}  Suppose
\(\operatorname{Hf}(\sigma)\) is valid but some \(\Lang\)-structure
\(\mathcal N\) satisfies \(\neg\sigma\), hence \(\sigma^{\neg}\).  By
Theorem~\ref{thm:skolem-equisatisfiable}(1), \(\mathcal N\) has an expansion
satisfying \(\operatorname{Sk}(\sigma^{\neg})=\forall\bar u\,\neg\theta^{\mathrm H}\),
which falsifies \(\operatorname{Hf}(\sigma)\).  This is a contradiction.
\end{proof}
\end{annexitem}

\begin{annexitem}{annex:p02-w02-p02-028-general-validity}{thm:p02-w02-p02-028-general-validity}{Herbrand's theorem, validity form for arbitrary sentences}
\begin{proof}
By Theorem~\ref{thm:prenex-normal-form}, \(\sigma\) and \(\sigma'\) are
logically equivalent and in the same signature, so one is valid iff the
other is.  The three operations composed in that theorem (negation normal
form, rectification, prenex extraction) rearrange connectives and
quantifiers and rename bound variables; they introduce no atom of the form
\(r=t\), so \(\sigma'\) is equality-free.  \(\Lang^{\mathrm H}\) contains
the constant of \(\Lang\).

If \(\sigma\) is valid, then \(\sigma'\) is valid, \(\operatorname{Hf}(\sigma')\)
is valid by the first implication of
Theorem~\ref{thm:p02-w02-p02-028-herbrand-form-validity} (no choice), and
Theorem~\ref{thm:p02-w02-p02-028-existential-validity}, applied to the
existential \(\Lang^{\mathrm H}\)-sentence \(\operatorname{Hf}(\sigma')\),
yields a propositionally valid finite expansion.  Conversely, if such an
expansion is propositionally valid, the right-to-left direction of
Theorem~\ref{thm:p02-w02-p02-028-existential-validity} makes
\(\operatorname{Hf}(\sigma')\) valid, the second implication of
Theorem~\ref{thm:p02-w02-p02-028-herbrand-form-validity} (with choice) makes
\(\sigma'\) valid, and so \(\sigma\) is valid.
\end{proof}
\end{annexitem}

\begin{annexitem}{annex:p02-w02-p02-028-clause-form}{thm:p02-w02-p02-028-clause-form}{Herbrand's theorem for clauses}
\begin{proof}
A ground clause \(K\) is propositionally satisfied by \(v\) when some literal
\(L\in K\) has \(\llbracket L\rrbracket_v=1\).

\emph{Right to left.}  Let \(G\) be a finite propositionally unsatisfiable
set of ground instances \(C\gamma\) of members of \(S\), and suppose
\(\mathcal N\) is a model of \(S\).  By Lemma~\ref{lem:clause-substitution},
\(\mathcal N\) satisfies every \(C\gamma\in G\); as \(C\gamma\) is ground,
some literal \(L\in C\gamma\) is true in \(\mathcal N\), so
\(\llbracket L\rrbracket_{v_{\mathcal N}}=1\) by
Lemma~\ref{thm:p02-w02-p02-028-ground-evaluation}(1).  Thus
\(v_{\mathcal N}\) satisfies \(G\), a contradiction.  So \(S\) is
unsatisfiable.

\emph{Left to right.}  Let \(\Lang_S\) consist of the symbols occurring in
\(S\) and, if none of them is a constant, one constant of \(\Lang\); it is
finite.  Suppose every finite set of ground \(\Lang_S\)-instances of members
of \(S\) is propositionally satisfiable; we show \(S\) satisfiable.  For
\(C\in S\) let \(\varphi_C\) be the disjunction of its literals in a fixed
order (\(\bottom\) for the empty clause); its variables are those of \(C\).
For a ground substitution \(\gamma\) on these variables,
\(\llbracket\varphi_C\gamma\rrbracket_v=1\) iff some \(L\in C\) has
\(\llbracket L\gamma\rrbracket_v=1\), iff \(v\) satisfies the ground clause
\(C\gamma\); identifying two literals in \(C\gamma\) does not change this.
Hence a finite set of ground instances \(\varphi_{C_i}\gamma_i\) is
propositionally satisfiable exactly when the clauses \(C_i\gamma_i\) are, and
the hypothesis of Lemma~\ref{thm:p02-w02-p02-028-herbrand-model} holds for
\(\Phi=\{\varphi_C:C\in S\}\) over \(\Lang_S\).  That lemma gives a Herbrand
\(\Lang_S\)-structure satisfying the universal closure of every
\(\varphi_C\), that is, a model of \(S\).  By the reduct fact it expands to an
\(\Lang\)-model of \(S\).  Contrapositively, if \(S\) is unsatisfiable, some
finite set of ground \(\Lang_S\)-instances, which are ground
\(\Lang\)-instances, is propositionally unsatisfiable.
\end{proof}
\end{annexitem}

\begin{annexitem}{annex:p02-w02-p02-028-validity-to-clauses}{thm:p02-w02-p02-028-validity-to-clauses}{Validity form and clause form}
\begin{proof}
Let \(\theta'=\operatorname{nnf}(\neg\theta^{\mathrm H})\).  By
Lemma~\ref{lem:fol-nnf} it is a quantifier-free formula in negation normal
form equivalent to \(\neg\theta^{\mathrm H}\), with the same atoms, and
\(S_\sigma=\operatorname{Cl}(\theta')\) is a finite equality-free clause
family whose variables are among \(\bar u\)
(Method~\ref{def:p02-matrix-clause-construction}).

\emph{(1)}  Fix a ground tuple \(\bar t\) and an atom valuation \(v\) for
ground \(\Lang^{\mathrm H}\)-atoms.  Let \(B=\{A:v(A)=1\}\), so
\(v_{\mathcal M_B}=v\), and let \(s\) be an assignment in \(\mathcal M_B\)
with \(s(u_i)=t_i\); as \(t_i^{\mathcal M_B}=t_i\), Lemma~\ref{thm:p02-w02-p02-028-ground-evaluation}(2)
applies to \(s\).  Then
\[
\begin{aligned}
\llbracket\neg\theta^{\mathrm H}(\bar t)\rrbracket_v=1
 &\iff \mathcal M_B\entails\neg\theta^{\mathrm H}(\bar t)
 \iff \mathcal M_B,s\entails\neg\theta^{\mathrm H}
 \iff \mathcal M_B,s\entails\theta'\\
 &\iff \text{every }C\in S_\sigma\text{ has some }L\in C\text{ with }\mathcal M_B,s\entails L\\
 &\iff \text{every }C\in S_\sigma\text{ has some }L\in C\text{ with }\llbracket L(\bar t)\rrbracket_v=1,
\end{aligned}
\]
using in turn Lemma~\ref{thm:p02-w02-p02-028-ground-evaluation}(1) and~(3),
part~(2) of the same lemma, Lemma~\ref{lem:fol-nnf},
Lemma~\ref{thm:p02-matrix-clause-equivalence}, and parts~(2), (1) and~(3)
again for each literal.  So \(v\) gives \(\neg\theta^{\mathrm H}(\bar t)\)
the value \(1\) iff \(v\) satisfies every ground clause \(C(\bar t)\).
Consequently the disjunction
\(\theta^{\mathrm H}(\bar t_1)\lor\cdots\lor\theta^{\mathrm H}(\bar t_m)\)
fails to be propositionally valid iff some \(v\) gives every
\(\neg\theta^{\mathrm H}(\bar t_j)\) the value \(1\), iff some \(v\)
satisfies \(\{C(\bar t_j):C\in S_\sigma,\ j\leq m\}\).

\emph{(2)}  Let \(C\gamma\) be a ground instance, \(\gamma\) sending the
variables of \(C\) to ground terms.  Fix a constant \(c\) of
\(\Lang^{\mathrm H}\) and let \(\bar t_\gamma\) have \(i\)-th member
\(\gamma(u_i)\) if \(u_i\) occurs in \(C\) and \(c\) otherwise.  Then
\(C(\bar t_\gamma)=C\gamma\).  A finite set of instances is therefore
contained in \(\{C(\bar t_\gamma):C\in S_\sigma,\ \gamma\text{ among those
used}\}\).

\emph{(3)}  Assume \(\sigma\) is valid.  By
Theorem~\ref{thm:p02-w02-p02-028-general-validity} (the direction without
choice) some expansion is propositionally valid; by item~(1) the
corresponding set of ground instances of \(S_\sigma\) is propositionally
unsatisfiable; by Theorem~\ref{thm:p02-w02-p02-028-clause-form} \(S_\sigma\)
is unsatisfiable.  Conversely, if \(S_\sigma\) is unsatisfiable,
Theorem~\ref{thm:p02-w02-p02-028-clause-form} gives a finite
propositionally unsatisfiable set of ground instances; by item~(2) it is
contained in a set of the form in item~(1), which is then unsatisfiable as
well; item~(1) makes the corresponding expansion propositionally valid; and
Theorem~\ref{thm:p02-w02-p02-028-general-validity} (with choice) makes
\(\sigma\) valid.
\end{proof}
\end{annexitem}

\begin{annexitem}{annex:p02-w02-p02-028-unification}{thm:p02-w02-p02-028-unification}{Termination and correctness of unification}
\begin{proof}
Substitutions compose as in \cref{sec:fol-cardinality-normalforms-04}, so
\(t(\alpha\beta)=(t\alpha)\beta\) and composition is associative.  A
substitution \(\eta\) \emph{solves} a finite set \(E\) of equations when
\(r\eta\) and \(t\eta\) are the same term for every \(r\doteq t\) in \(E\); every
substitution solves the empty set.  A \emph{state} of a run is the pair
\((E,\theta)\) of the current equation set and the accumulator.  The size
\(\abs t\) of a term is its number of symbol occurrences (variables,
constants and function symbols).

\emph{Two facts about terms.}  (F1) If \(x\) does not occur in \(t\) and
\(\delta=\{x\mapsto t\}\), then \(t\delta=t\); more generally
\(r\delta=r\) for every term \(r\) without \(x\) (induction on \(r\)).  (F2) If
\(x\) occurs in \(t\) and \(t\) is not the variable \(x\), then for every
\(\eta\), \(\abs{t\eta}>\abs{x\eta}\), so \(x\eta\neq t\eta\).  Indeed \(t\) is
then \(f(t_1,\ldots,t_k)\) with \(x\) occurring in some \(t_i\); by induction on
terms, \(\abs{t_i\eta}\geq\abs{x\eta}\) whenever \(x\) occurs in \(t_i\)
(equality when \(t_i\) is \(x\)), and
\(\abs{t\eta}=1+\sum_j\abs{t_j\eta}>\abs{t_i\eta}\).

\emph{The procedure never gets stuck.}  Let \(r\doteq t\) be any equation.
If \(r\) is \(t\), deletion applies.  If neither is a variable, either their
outer symbols and arities agree and decomposition applies, or the clash rule
applies.  If \(r\) is not a variable and \(t\) is one, orientation applies.  If
\(r\) is a variable different from \(t\), the occurs check either fails or
elimination applies.  So a run stops only by failure or with \(E\) empty.

\emph{(1) Termination.}  To a finite set \(E\) assign the triple
\(\mu(E)=(n_1,n_2,n_3)\), where \(n_1\) is the number of distinct variables
occurring in \(E\), \(n_2=\sum_{(r\doteq t)\in E}(\abs r+\abs t)\), and
\(n_3\) is the number of equations of \(E\) whose left side is not a
variable and whose right side is a variable.  Triples of natural numbers
under the lexicographic order are well-ordered, so it suffices that every
transformation that does not fail strictly decreases \(\mu\).  Since \(E\) is a
set, adding an equation already present adds nothing.
\begin{itemize}
  \item \emph{Deletion} of \(t\doteq t\): no variable is added, so \(n_1\)
        does not increase, and \(n_2\) decreases by \(2\abs t\geq2\).
  \item \emph{Decomposition} of \(f(r_1,\ldots,r_n)\doteq f(t_1,\ldots,t_n)\)
        (\(n\geq0\)): the new equations contain exactly the variables of the
        removed one, so \(n_1\) does not increase; the removed equation
        contributes \(2+\sum_i(\abs{r_i}+\abs{t_i})\) to \(n_2\) and the
        added ones at most \(\sum_i(\abs{r_i}+\abs{t_i})\), so \(n_2\)
        decreases by at least \(2\).
  \item \emph{Orientation} of \(t\doteq x\) to \(x\doteq t\), with \(t\) not a
        variable: the variables are unchanged, \(n_2\) does not increase,
        and \(n_3\) decreases by one if \(n_2\) stays the same.
  \item \emph{Elimination} of \(x\doteq t\), with \(x\) not in \(t\): the new
        set is \(E'=\{r\delta\doteq t'\delta:(r\doteq t')\in E,\ (r\doteq t')\neq(x\doteq t)\}\)
        with \(\delta=\{x\mapsto t\}\).  Each \(r\delta\) has its variables
        among those of \(r\) other than \(x\), together with those of \(t\);
        all of these occur in \(E\), and \(x\) occurs in none of them.  So
        the variables of \(E'\) are among those of \(E\) other than \(x\),
        while \(x\) occurs in \(E\); hence \(n_1\) decreases.
\end{itemize}

\emph{The invariant.}  We show that every state \((E,\theta)\) reached
satisfies
\[
 (\mathrm{I})\qquad
 \begin{gathered}
 \text{the unifiers of }A\text{ and }B\text{ are exactly}\\
 \text{the substitutions }\theta\eta\text{ with }\eta\text{ solving }E .
 \end{gathered}
\]
\emph{Initial state.}  If the outer relation symbols or arities of \(A\) and
\(B\) differ, there is no unifier, since a substitution changes neither, and
the run fails correctly at once.  Otherwise
\(A=R(r_1,\ldots,r_k)\), \(B=R(t_1,\ldots,t_k)\), \(E=\{r_i\doteq t_i\}\) and
\(\theta=\mathrm{id}\); \(A\eta\) and \(B\eta\) are the same atom iff
\(r_i\eta\) is \(t_i\eta\) for each \(i\), that is iff \(\eta\) solves \(E\),
and \(\mathrm{id}\,\eta=\eta\).  So (I) holds.

\emph{Preservation.}  Let \((E,\theta)\) satisfy (I).
\begin{itemize}
  \item \emph{Deletion, decomposition, orientation} leave \(\theta\) unchanged
        and do not change the set of solutions: every \(\eta\) solves
        \(t\doteq t\); \(f(\bar r)\eta\) is \(f(\bar t)\eta\) iff
        \(r_i\eta\) is \(t_i\eta\) for each \(i\), because two terms with the
        same outer function symbol are identical iff their arguments are; and
        \(t\eta\) is \(x\eta\) iff \(x\eta\) is \(t\eta\).  So (I) persists.
  \item \emph{Elimination} of \(x\doteq t\) produces \((E',\theta\delta)\).
        First, if \(x\eta\) is \(t\eta\), then \(\delta\eta=\eta\): the two
        agree at \(x\), where \((\delta\eta)(x)=t\eta\), and at every other
        variable \(y\), where \((\delta\eta)(y)=y\eta\).  Suppose \(\eta\)
        solves \(E\).  Then \(\delta\eta=\eta\), and for each other equation
        \(r\doteq t'\) of \(E\), \((r\delta)\eta=r\eta=t'\eta=(t'\delta)\eta\),
        so \(\eta\) solves \(E'\) and \(\theta\eta=(\theta\delta)\eta\).
        Conversely, suppose \(\eta'\) solves \(E'\) and put
        \(\eta=\delta\eta'\).  Then \(x\eta=t\eta'\) and, by (F1),
        \(t\eta=(t\delta)\eta'=t\eta'\), so \(\eta\) solves \(x\doteq t\);
        for each other equation, \(r\eta=(r\delta)\eta'=(t'\delta)\eta'=t'\eta\).
        So \(\eta\) solves \(E\), and \((\theta\delta)\eta'=\theta\eta\).
        Hence the substitutions \(\theta\eta\) with \(\eta\) solving \(E\) are
        exactly the substitutions \((\theta\delta)\eta'\) with \(\eta'\)
        solving \(E'\), and (I) persists.
\end{itemize}

\emph{(2) and (3).}  By (1) and the fact that the procedure never gets
stuck, the run ends either in failure or with \(E=\varnothing\).  At the clash
rule, no \(\eta\) solves the selected equation, because
\(f(\bar r)\eta\) and \(g(\bar t)\eta\) have different outer symbols or
arities; at the occurs check, no \(\eta\) solves it, by (F2).  In both cases no
\(\eta\) solves \(E\), and (I) says that \(A\) and \(B\) have no unifier.  If
the run ends with \(E=\varnothing\) and accumulator \(\theta\), every
substitution solves \(E\), so by (I) the unifiers are exactly the
substitutions \(\theta\eta\).  Taking \(\eta=\mathrm{id}\), \(\theta\) is a
unifier, and every unifier has the form \(\theta\eta\); so \(\theta\) is a most
general unifier.  Consequently a run fails iff there is no unifier, and two
unifiable atoms have a most general unifier.
\end{proof}
\end{annexitem}

\begin{annexitem}{annex:p02-w02-p02-028-lifting}{thm:p02-w02-p02-028-lifting}{Lifting of ground resolution}
\begin{proof}
For a ground atom \(P\) write \(P^+=P\) and \(P^-=\neg P\).  All
substitutions compose as in \cref{sec:fol-cardinality-normalforms-04}, so
\(X(\alpha\beta)=(X\alpha)\beta\) for terms, literals and clauses.

\emph{Claim A (factoring).}  Let \(E\) occur in a first-order derivation from
\(S\), let \(\gamma\) send the variables of \(E\) to ground terms, let \(P\)
be a ground atom and \(\lambda\in\{+,-\}\), and let
\(F=\{L\in E:L\gamma=P^\lambda\}\).  Then the derivation can be extended by
Factoring steps to contain \(E\alpha\) for some substitution \(\alpha\), and
there is \(\gamma'\) with \((L\alpha)\gamma'=L\gamma\) for every \(L\in E\)
and with \(F\alpha\) having at most one member.

Start with \(\alpha=\mathrm{id}\), \(\gamma'=\gamma\).  If \(F\alpha\) has two
distinct members \(L_1,L_2\), they have polarity \(\lambda\), and their atoms
\(A_1,A_2\) satisfy \(A_1\gamma'=A_2\gamma'=P\), because each \(L_k\) is
\(L\alpha\) for some \(L\in F\).  So \(\gamma'\) unifies \(A_1\) and
\(A_2\); by Lemma~\ref{thm:p02-w02-p02-028-unification} they have a most
general unifier
\(\mu\), and Factoring derives \((E\alpha)\mu=E(\alpha\mu)\).  By the
definition of most general unifier, \(\gamma'=\mu\eta\) for some \(\eta\).
Replace \(\alpha\) by \(\alpha\mu\) and \(\gamma'\) by \(\eta\); then
\((L\alpha\mu)\eta=(L\alpha)\gamma'=L\gamma\).  Since \(L_1\mu=L_2\mu\), the
set \(F\alpha\mu\) is strictly smaller than \(F\alpha\).  The iteration
therefore stops with \(F\alpha\) of size at most one.  Moreover
\((E\alpha)\gamma'=E\gamma\) is ground, so \(\gamma'\) sends every variable
of \(E\alpha\) to a ground term.

\emph{Claim B (one resolution step).}  Let \(E,F\) occur in a derivation from
\(S\), let \(\gamma_E,\gamma_F\) be ground on their variables, and let
\(G=(E\gamma_E\setminus\{P\})\cup(F\gamma_F\setminus\{\neg P\})\) with
\(P\in E\gamma_E\) and \(\neg P\in F\gamma_F\).  Then the derivation extends
by Factoring and one Resolution step to contain a clause \(R\), and some
\(\eta\) sends the variables of \(R\) to ground terms with \(R\eta=G\).

Apply Claim~A to \(E,\gamma_E,P,+\) and to \(F,\gamma_F,P,-\), obtaining
\(E^*=E\alpha\) with \(\gamma_E^*\) and \(F^*=F\beta\) with
\(\gamma_F^*\).  Since \(P\in E\gamma_E\), the set \(\{L\in E:L\gamma_E=P\}\)
is nonempty, so its image is a single positive literal \(A^*\in E^*\) with
\(A^*\gamma_E^*=P\).  It is the only literal of \(E^*\) sent to \(P\): if
\(L\alpha\in E^*\) has \((L\alpha)\gamma_E^*=P\), then \(L\gamma_E=P\), so
\(L\alpha=A^*\).  Likewise \(\neg B^*\in F^*\) is the only literal of \(F^*\)
sent to \(\neg P\) by \(\gamma_F^*\).  Choose renamings \(\pi,\rho\) making
\(E^*\pi\) and \(F^*\rho\) renamed apart.  Define \(\delta\) by
\(\delta(x)=(\pi^{-1}x)\gamma_E^*\) for \(x\) occurring in \(E^*\pi\),
\(\delta(x)=(\rho^{-1}x)\gamma_F^*\) for \(x\) occurring in \(F^*\rho\), and
\(\delta(x)=x\) otherwise; it has finite support, and
\((L\pi)\delta=L\gamma_E^*\) for \(L\in E^*\),
\((L\rho)\delta=L\gamma_F^*\) for \(L\in F^*\).  Then
\((A^*\pi)\delta=P=(B^*\rho)\delta\), so \(A^*\pi\) and \(B^*\rho\) are
unifiable; by Lemma~\ref{thm:p02-w02-p02-028-unification} they have a most
general unifier \(\nu\), and \(\delta=\nu\eta\) for some \(\eta\).
Resolution on \(A^*\) and \(\neg B^*\) gives
\(R=\bigl((E^*\pi\setminus\{A^*\pi\})\cup(F^*\rho\setminus\{\neg B^*\rho\})\bigr)\nu\),
and
\[
 R\eta=(E^*\pi\setminus\{A^*\pi\})\delta\cup(F^*\rho\setminus\{\neg B^*\rho\})\delta .
\]
Now \((E^*\pi)\delta=E^*\gamma_E^*=E\gamma_E\), and \(A^*\pi\) is the only
member of \(E^*\pi\) sent to \(P\) (renaming is injective on literals).
Hence \((E^*\pi\setminus\{A^*\pi\})\delta=E\gamma_E\setminus\{P\}\): every
other member is sent to a member of \(E\gamma_E\) different from \(P\), and
every member of \(E\gamma_E\) different from \(P\) is the image of a member
of \(E^*\pi\) other than \(A^*\pi\).  The same holds for \(F\), so
\(R\eta=G\).  As \(G\) is ground, \(\eta\) sends every variable of \(R\) to a
ground term.

\emph{Induction.}  We define \(K_j'\) and \(\gamma_j\) by recursion on \(j\)
and keep one first-order derivation from \(S\) containing all clauses defined
so far.  If \(K_j=C\gamma\) with \(C\in S\), add \(C\) and put
\(K_j'=C\), \(\gamma_j=\gamma\).  If \(K_j\) is the propositional resolvent of
\(K_a\) and \(K_b\) (\(a,b<j\)) on a ground atom \(P\), let \(E\) be the one
of \(K_a',K_b'\) whose instance contains \(P\) and \(F\) the other; Claim~B
with \(\gamma_E,\gamma_F\) the corresponding \(\gamma_a,\gamma_b\) gives
\(K_j'=R\) and \(\gamma_j=\eta\) with \(K_j'\gamma_j=K_j\).  If
\(K_r=\varnothing\), then \(K_r'\gamma_r=\varnothing\), which forces
\(K_r'=\varnothing\), since a substitution sends a nonempty clause to a
nonempty clause.
\end{proof}
\end{annexitem}

\begin{annexitem}{annex:p02-w02-p02-028-resolution}{thm:p02-w02-p02-028-resolution}{From the clause form to resolution}
\begin{proof}
(1) implies (2): by Theorem~\ref{thm:p02-w02-p02-028-clause-form} some finite
set \(G\) of ground instances is propositionally unsatisfiable.  Its finitely
many ground atoms can be renamed injectively by propositional variables
\(p_n\); the renaming commutes with forming resolvents and preserves
satisfiability, so Theorem~\ref{thm:prop-resolution-complete} gives a
resolution refutation of the renamed set, and renaming back gives one of
\(G\).  (2) implies (3) by Lemma~\ref{thm:p02-w02-p02-028-lifting}.  (3)
implies (1) by the soundness theorem for first-order Resolution and Factoring
in \cref{sec:fol-cardinality-normalforms-04}.  The final assertion combines
this equivalence for \(S_\sigma\) with
Theorem~\ref{thm:p02-w02-p02-028-validity-to-clauses}(3).
\end{proof}
\end{annexitem}

\begin{annexitem}{annex:p02-w02-p02-028-no-recursive-bound}{thm:p02-w02-p02-028-no-recursive-bound}{Search for instances and the absence of a recursive bound}
\begin{proof}
(1)  For a given existential sentence, the ground terms over its symbols
(with one constant added if it has none) can be effectively listed, hence so
can all finite lists of tuples.  For each list, propositional validity of the
expansion is decided by a truth table on its finitely many atoms.  The
procedure tests the lists in order and halts at the first success.  By the
right-to-left direction of
Theorem~\ref{thm:p02-w02-p02-028-existential-validity} it halts only on valid
sentences, and by the left-to-right direction it halts on every valid one.
If a derivation of the sentence in some sound calculus with effectively
recognizable derivations is supplied, the same search, applied to the
conclusion of the derivation, is a total recursive function of the
derivation, because soundness makes every derived sentence valid; this is the
first observation of [1.2] in the equality-free case.

(2)  Suppose such a total recursive function \(F\) existed.  Given an
equality-free existential sentence \(\tau\), compute \(F(\tau)\) and test the
expansion by truth table.  If it is valid, \(\tau\) is valid by the
right-to-left direction of
Theorem~\ref{thm:p02-w02-p02-028-existential-validity}; if not, \(\tau\) is
not valid, since otherwise \(F(\tau)\) would give a valid expansion.  So
validity of equality-free existential sentences would be decidable.  For an
arbitrary equality-free sentence \(\sigma\), the prenex form \(\sigma'\) of
Theorem~\ref{thm:prenex-normal-form} and the Herbrand form
\(\operatorname{Hf}(\sigma')\) are computed by explicit recursions, and
\(\sigma\) is valid iff \(\operatorname{Hf}(\sigma')\) is
(Theorem~\ref{thm:p02-w02-p02-028-herbrand-form-validity}).  Hence
equality-free validity would be decidable, against the hypothesis.
\end{proof}
\end{annexitem}

\subsection*{Statman's chapter}

In the remaining items, terms, computations and interpretations are those of
Section~\ref{sec:p02-w02-p02-028-statman}.  Two facts are used repeatedly.
\emph{(T1) Term values.}  Under any interpretation \(I\) in the sense of
Convention~\ref{def:p02-w02-p02-028-st-convention}, if \(I(a)=I(b)\) then
\(I(d(a/x))=I(d(b/x))\) for every term \(d\); this is an induction on \(d\),
the case \(fd_1\cdots d_n\) using that \(I(f)\) is a function.  Hence rule
(1) preserves truth under \(I\), rule (2) does because successor is
injective, and an axiom \(a=a\) is true.  \emph{(T2) Tree size.}  A
computation whose branches contain at most \(h\) equation occurrences has
length at most \(2^h-1\), since each node has at most two premises.

The first two items concern first-order logic with equality.  In them
\(\approx\) is a binary relation symbol not in \(\Lang\), and for a formula
\(\varphi\) of \(\Lang\), \(\varphi^{\approx}\) replaces every atom \(r=t\) by
\(r\approx t\); so \(\varphi^{\approx}\) is equality-free.  We use the
following form of Lemma~\ref{thm:p02-w02-p02-028-ground-evaluation}(2), which
its proof establishes, since that proof shows
\(\sem{r(\bar t)}^{\mathcal N}_s=\sem{r}^{\mathcal N}_{s'}\) for every
assignment \(s\) and then argues atom by atom and connective by connective:
\emph{(G) for a quantifier-free \(\varphi\), ground terms \(t_1,\ldots,t_k\),
distinct variables \(z_1,\ldots,z_k\) and an assignment \(s\) with
\(s(z_i)=t_i^{\mathcal N}\), \(\mathcal N,s\entails\varphi\) iff
\(\mathcal N,s\entails\varphi(t_1/z_1,\ldots,t_k/z_k)\).}

\begin{annexitem}{annex:p02-w02-p02-028-st-herbrand-eq}{thm:p02-w02-p02-028-st-herbrand-eq}{Herbrand's theorem with equality}
\begin{proof}
\emph{Converse.}  If \(\bigvee_iG(\bar t^i)\) is true in \(\mathcal N\) under
\(s\), some disjunct \(G(\bar t^i)\) is.  The proof of
Lemma~\ref{thm:p02-w02-p02-028-ground-evaluation}(2) goes through for
arbitrary terms \(t^i_j\) when \(s'\) is taken to be
\(s[x_j\mapsto\sem{t^i_j}^{\mathcal N}_s]\): in the variable case the left
side is \(\sem{t^i_j}^{\mathcal N}_s=s'(x_j)\), and no other step uses that
the terms are ground, since \(G\) has no quantifiers.  So
\(\mathcal N,s'\entails G\), and \(\mathcal N,s\entails\exists\bar x\,G\).

\emph{Main direction.}  Let \(y_1,\ldots,y_l\) be the free variables of \(G\)
other than \(x_1,\ldots,x_n\).  Let \(d_1,\ldots,d_l\) and \(c_0\) be
constants not in \(\Lang\), and \(G_0=G(d_1/y_1,\ldots,d_l/y_l)\).  Let
\(\Lang_0\) be the finite signature of the symbols occurring in \(G_0\)
together with \(c_0\).

\emph{Step 1: \(\exists\bar x\,G_0\) is true in every
\(\Lang_0\)-structure.}  An \(\Lang_0\)-structure \(\mathcal N\) expands to
an \(\Lang\cup\{d_1,\ldots,d_l,c_0\}\)-structure \(\mathcal N^+\) as in the
reduct fact.  Take \(s\) with \(s(y_k)=d_k^{\mathcal N^+}\).  By validity some
tuple \(\bar a\) has \(\mathcal N^+,s[\bar x\mapsto\bar a]\entails G\), hence
\(\mathcal N^+,s[\bar x\mapsto\bar a]\entails G_0\) by (G).  So
\(\exists\bar x\,G_0\) holds in \(\mathcal N^+\), and in \(\mathcal N\) by
the reduct fact.

\emph{Step 2: the equality axioms.}  Let \(\mathrm{Eq}\) consist of the
following quantifier-free formulas over fresh variables, distinct from each
other and from \(\bar x\): \(z\approx z\); \(z\approx z'\to z'\approx z\);
\((z\approx z'\land z'\approx z'')\to z\approx z''\); for each \(k\)-ary
function symbol \(f\) of \(\Lang_0\),
\((z_1\approx z_1'\land\cdots\land z_k\approx z_k')\to f(\bar z)\approx f(\bar z')\);
and for each \(k\)-ary relation symbol \(R\) of \(\Lang_0\),
\((z_1\approx z_1'\land\cdots\land z_k\approx z_k'\land R(\bar z))\to R(\bar z')\).
Rename variables so that distinct members share none, let \(\varepsilon\) be
their conjunction and \(\bar w\) its variables.  Put
\(\tau=\exists\bar w\,\exists\bar x\,(\neg\varepsilon\lor G_0^{\approx})\), an
equality-free existential sentence of \(\Lang_0\cup\{\approx\}\).

\emph{Step 3: \(\tau\) is valid.}  Let \(\mathcal N'\) be an
\(\Lang_0\cup\{\approx\}\)-structure, with \(\approx\) read as an arbitrary
binary relation.  If some values of \(\bar w\) falsify \(\varepsilon\), those
values and any values of \(\bar x\) witness \(\tau\).  Otherwise every member
of \(\mathrm{Eq}\) holds under every assignment, so the relation
\(\sim\) interpreting \(\approx\) is an equivalence relation, and
\(a_i\sim a_i'\) for all \(i\) implies
\(f^{\mathcal N'}(\bar a)\sim f^{\mathcal N'}(\bar a')\) and
\(R^{\mathcal N'}(\bar a)\Rightarrow R^{\mathcal N'}(\bar a')\).  Let
\(\mathcal Q\) have domain the set of classes \([a]\), and put
\(f^{\mathcal Q}([a_1],\ldots,[a_k])=[f^{\mathcal N'}(a_1,\ldots,a_k)]\),
\(c^{\mathcal Q}=[c^{\mathcal N'}]\), and
\(R^{\mathcal Q}([a_1],\ldots,[a_k])\) iff \(R^{\mathcal N'}(a_1,\ldots,a_k)\);
these are independent of the representatives by the two compatibility
properties (with symmetry for \(R\)).  Equality in \(\mathcal Q\) is identity.
For an assignment \(s\) in \(\mathcal N'\) let \([s](v)=[s(v)]\).  Induction
on terms gives \(\sem t^{\mathcal Q}_{[s]}=[\sem t^{\mathcal N'}_s]\).  Hence
for an atom, \(\mathcal Q,[s]\entails r=t\) iff
\(\sem r^{\mathcal N'}_s\sim\sem t^{\mathcal N'}_s\) iff
\(\mathcal N',s\entails r\approx t\), and \(\mathcal Q,[s]\entails R(\bar t)\)
iff \(\mathcal N',s\entails R(\bar t)\); by induction on connectives,
\(\mathcal Q,[s]\entails\varphi\) iff \(\mathcal N',s\entails\varphi^{\approx}\)
for every quantifier-free \(\Lang_0\)-formula \(\varphi\).  Fix any
assignment \(s_0\) in \(\mathcal N'\).  By Step~1, \(\exists\bar x\,G_0\) is
true in the \(\Lang_0\)-structure \(\mathcal Q\), so there are classes
\([a_1],\ldots,[a_n]\) with
\(\mathcal Q,[s_0][\bar x\mapsto[\bar a]]\entails G_0\).  As
\([s_0][\bar x\mapsto[\bar a]]=[s_0[\bar x\mapsto\bar a]]\), we get
\(\mathcal N',s_0[\bar x\mapsto\bar a]\entails G_0^{\approx}\), and \(\tau\)
holds in \(\mathcal N'\) with any values for \(\bar w\).

\emph{Step 4: Herbrand instances.}  \(\Lang_0\cup\{\approx\}\) has the
constant \(c_0\).  By Theorem~\ref{thm:p02-w02-p02-028-existential-validity}
there are ground tuples \(\bar u^1,\bar t^1,\ldots,\bar u^m,\bar t^m\)
(\(m\geq1\)) with
\(D=\bigvee_{i=1}^m\bigl(\neg\varepsilon(\bar u^i)\lor G_0^{\approx}(\bar t^i)\bigr)\)
propositionally valid.

\emph{Step 5: discarding the axioms.}  Let \(\mathcal N\) be any structure for
\(\Lang\cup\{d_1,\ldots,d_l,c_0\}\), and let \(\mathcal N^{=}\) expand it by
interpreting \(\approx\) as identity.  Every ground instance of a member of
\(\mathrm{Eq}\) is true in \(\mathcal N^{=}\), because identity is an
equivalence relation respected by every function and relation; so each
\(\varepsilon(\bar u^i)\) is true.  \(D\) is true in \(\mathcal N^{=}\) by
Lemma~\ref{thm:p02-w02-p02-028-ground-evaluation}(1).  Hence some
\(G_0^{\approx}(\bar t^i)\) is true in \(\mathcal N^{=}\), and since \(\approx\)
is identity there, \(G_0(\bar t^i)\) is true in \(\mathcal N\).  Thus
\(\bigvee_iG_0(\bar t^i)\) is valid.

\emph{Step 6: returning to \(\Lang\).}  Let \(z\) be a variable not occurring
in \(G\), and let \(t^i_j\) be the \(\Lang\)-term obtained from the
\(j\)-th member of \(\bar t^i\) by replacing each \(d_k\) by \(y_k\) and
\(c_0\) by \(z\).  Replacing, in \(\bigvee_iG(t^i_1,\ldots,t^i_n)\), each
\(y_k\) by \(d_k\) and \(z\) by \(c_0\) gives \(\bigvee_iG_0(\bar t^i)\),
because the \(x_j\) are replaced first and \(z\) does not occur in \(G\).  Given
an \(\Lang\)-structure \(\mathcal N\) and an assignment \(s\), expand
\(\mathcal N\) by \(d_k\mapsto s(y_k)\) and \(c_0\mapsto s(z)\).  The
expansion satisfies \(\bigvee_iG_0(\bar t^i)\) by Step~5, so by (G) it
satisfies \(\bigvee_iG(t^i_1,\ldots,t^i_n)\) under \(s\), and so does
\(\mathcal N\) by the reduct fact.  Hence the disjunction is valid.
\end{proof}
\end{annexitem}

\begin{annexitem}{annex:p02-w02-p02-028-st-methods}{thm:p02-w02-p02-028-st-methods}{Finding the terms}
\begin{proof}
(1)  Let \(\psi\) be quantifier-free.  Replace its free variables by distinct
new constants to obtain a ground \(\psi'\); by (G) and the reduct fact,
\(\psi\) is valid iff \(\psi'\) is (to pass from \(\psi'\) back to \(\psi\),
interpret each new constant by the value of its variable).  Let \(U\) be the
finite set of subterms of terms occurring in \(\psi'\), together with one
new constant \(c_*\).  Let \(\Gamma\) be the finite set of the following
ground formulas: \(t\approx t\), \(t\approx t'\to t'\approx t\) and
\((t\approx t'\land t'\approx t'')\to t\approx t''\) for \(t,t',t''\in U\);
\((\bigwedge_ia_i\approx b_i)\to f(\bar a)\approx f(\bar b)\) whenever
\(f(\bar a)\) and \(f(\bar b)\) both belong to \(U\); and
\((\bigwedge_ia_i\approx b_i\land R(\bar a))\to R(\bar b)\) for every relation
symbol \(R\) of \(\psi'\) and all tuples \(\bar a,\bar b\) from \(U\).  We
claim that \(\psi'\) is valid iff \(\bigwedge\Gamma\to\psi'^{\approx}\) is
propositionally valid; the latter is decided by a truth table, which proves
(1).

If \(\bigwedge\Gamma\to\psi'^{\approx}\) is propositionally valid and
\(\mathcal N\) is any structure, expand it by reading \(\approx\) as
identity; every member of \(\Gamma\) is then true, and so is the implication
(Lemma~\ref{thm:p02-w02-p02-028-ground-evaluation}(1)), so
\(\psi'^{\approx}\) is true, that is, \(\psi'\) is true in \(\mathcal N\).
Conversely, let an atom valuation \(v\) make every member of \(\Gamma\) true
and \(\psi'^{\approx}\) false.  Put \(t\sim t'\) iff \(v(t\approx t')=1\), an
equivalence relation on \(U\).  Let \(\mathcal Q\) have domain the classes of
\(\sim\); interpret a constant \(c\in U\) by \([c]\); interpret a function
symbol \(f\) occurring in \(\psi'\) by
\(f^{\mathcal Q}([a_1],\ldots,[a_k])=[f(a_1',\ldots,a_k')]\) if some
\(f(a_1',\ldots,a_k')\in U\) has \(a_i'\sim a_i\) for all \(i\), and by
\([c_*]\) otherwise, which is well defined by the congruence members for
\(f\); interpret \(R\) by \(R^{\mathcal Q}([a_1],\ldots,[a_k])\) iff some
\(\bar a'\) from \(U\) with \(a_i'\sim a_i\) has \(v(R(\bar a'))=1\); and
interpret every other symbol arbitrarily.  By induction,
\(t^{\mathcal Q}=[t]\) for \(t\in U\), since \(f(\bar a)\in U\) is its own
witness.  For an atom \(r=t\) of \(\psi'\), \(\mathcal Q\entails r=t\) iff
\([r]=[t]\) iff \(v(r\approx t)=1\).  For an atom \(R(\bar a)\) of \(\psi'\),
\(\mathcal Q\entails R(\bar a)\) iff some \(\bar a'\sim\bar a\) has
\(v(R(\bar a'))=1\), iff \(v(R(\bar a))=1\), by the congruence members for
\(R\) and symmetry.  By induction on connectives, \(\psi'\) is true in
\(\mathcal Q\) iff \(v\) gives \(\psi'^{\approx}\) the value \(1\), which it
does not.  So \(\psi'\) is not valid.

(2)  Under the coding, the finite lists of \(n\)-tuples of terms can be
effectively listed.  For each list in turn, decide validity of
\(\bigvee_iG(\bar t^i)\) by (1); stop at the first valid one.  By the
converse in Theorem~\ref{thm:p02-w02-p02-028-st-herbrand-eq} the procedure
stops only on valid formulas, and by its main direction it stops on every
valid one.  After stopping, list the natural deduction derivations without
open assumptions and stop at the first whose conclusion is the disjunction;
by Theorem~\ref{thm:fol-completeness} (with the Axiom of Choice) one exists,
and checking that a coded object is such a derivation is effective.

(3)  By Theorem~\ref{thm:fol-soundness}, the conclusion of a derivation
without open assumptions is valid; apply the procedure of (2) to it.  It
stops for every input derivation, so the composite is total recursive.

(4)  Suppose \(F\) were such a total recursive procedure.  For a formula
\(\exists\bar x\,G\), compute \(F\)'s list and decide by (1) whether its
disjunction is valid.  If it is, \(\exists\bar x\,G\) is valid by the
converse in Theorem~\ref{thm:p02-w02-p02-028-st-herbrand-eq}; if not,
\(\exists\bar x\,G\) is not valid, since otherwise \(F\) would give a valid
disjunction.  So validity of such formulas would be decidable.  Given any
sentence \(\sigma\), compute a rectified prenex form \(\sigma'\)
(Theorem~\ref{thm:prenex-normal-form}) and its Herbrand form
\(\operatorname{Hf}(\sigma')\), choosing the new function symbols from those
not in \(\sigma\), as the hypothesis on the coding allows.  The proof of
Theorem~\ref{thm:p02-w02-p02-028-herbrand-form-validity} uses only
Lemma~\ref{thm:p02-w02-p02-028-herbrand-skolem-duality}, the quantifier
dualities and Theorem~\ref{thm:skolem-equisatisfiable}, none of which is
restricted to equality-free sentences, so with the Axiom of Choice it shows
that \(\sigma\) is valid iff \(\operatorname{Hf}(\sigma')\) is.  Hence
validity of all sentences would be decidable, contrary to
Theorem~\ref{thm:fol-validity-undecidable}.
\end{proof}
\end{annexitem}

\begin{annexitem}{annex:p02-w02-p02-028-st-measure-transfer}{thm:p02-w02-p02-028-st-measure-transfer}{Measures under non-projecting substitution}
\begin{proof}
Let \(c\), \(x\), \(a\), \(b\) be given.  Choose an individual variable
\(z\) that occurs in none of \(a,b,c\), is not in the domain of \(\theta\),
and occurs in no value \(\theta f\) or \(\theta y\).  Put \(c_0=c(z/x)\);
then \(c(a/x)=c_0(a/z)\) and \(c(b/x)=c_0(b/z)\).

\emph{Claim 1: \(\theta(c_0(e/z))=(\theta c_0)(\theta e/z)\) for every term
\(e\).}  Induction on \(c_0\).  If \(c_0\) is \(z\), both sides are \(\theta e\),
as \(\theta z=z\).  If \(c_0\) is another variable or \(0\), both sides are
\(\theta c_0\), which does not contain \(z\).  The case \(sd\) is immediate
from the induction hypothesis.  If \(c_0=fd_1\cdots d_n\) and \(f\) is not in
the domain of \(\theta\), use the induction hypothesis on each \(d_i\).  If
\(\theta f=\lambda y_1\cdots y_n\,g\), then
\[
 \theta(c_0(e/z))=g\bigl(\theta(d_1(e/z))/y_1,\ldots\bigr)
 =g\bigl((\theta d_1)(\theta e/z)/y_1,\ldots\bigr),
\]
while
\((\theta c_0)(\theta e/z)=\bigl(g(\theta d_1/y_1,\ldots)\bigr)(\theta e/z)\);
these agree because \(z\) does not occur in \(g\), so substituting
\(\theta e\) for \(z\) affects only the inserted terms \(\theta d_i\).

\emph{Claim 2: if \(z\) occurs in \(c_0\), it occurs in \(\theta c_0\).}
Induction on \(c_0\): the cases \(z\) and \(sd\) are clear.  For
\(fd_1\cdots d_n\) with \(z\) in \(d_i\): if \(f\) is not in the domain of
\(\theta\), \(\theta c_0\) contains \(\theta d_i\); if
\(\theta f=\lambda y_1\cdots y_n\,g\), then \(y_i\) occurs in \(g\) because
\(\theta\) is non-projecting, so \(g(\theta d_1/y_1,\ldots)\) contains
\(\theta d_i\), which contains \(z\) by the induction hypothesis.

Now, if \(M^*(a)\leq M^*(b)\), that is \(M(\theta a)\leq M(\theta b)\), then
by Claim~1 and the first measure condition for \(M\),
\(M^*(c(a/x))=M((\theta c_0)(\theta a/z))\leq M((\theta c_0)(\theta b/z))=M^*(c(b/x))\).
If \(x\) occurs in \(c\), then \(z\) occurs in \(c_0\) and, by Claim~2, in
\(\theta c_0\); the second measure condition for \(M\) gives
\(M^*(a)=M(\theta a)\leq M((\theta c_0)(\theta a/z))=M^*(c(a/x))\).  For the
counterexample in the statement, \(\theta(fx)=0\) has length \(1\) and
\(\theta(sx)=sx\) has length \(2\).
\end{proof}
\end{annexitem}

\begin{annexitem}{annex:p02-w02-p02-028-st-moves-basic}{thm:p02-w02-p02-028-st-moves-basic}{Moves preserve length, conclusion and singularity}
\begin{proof}
In each row of Figure~\ref{fig:p02-w02-p02-028-st-moves}, the two shapes have
the same root, the same three immediate subcomputations \(T_1,T_2,T_3\), and
two inferred equation occurrences each.  So a move preserves the conclusion,
the assumptions and axioms (the leaves), and the length
\(\operatorname{lh}(T_1)+\operatorname{lh}(T_2)+\operatorname{lh}(T_3)+2\);
induction on the number of moves extends this to reduction.

The new inferences are instances of rule (1).  Switch: from
\(E(a/x,b/y)\) and \(b\doteq d\) infer \(E(a/x,d/y)\), substituting at \(y\)
in \(E(a/x,\cdot/y)\); then from that and \(a\doteq c\) infer
\(E(c/x,d/y)\).  First shift: from \(a\doteq c(b/y)\) and \(b\doteq d\)
infer \(a\doteq c(d/y)\), substituting at \(y\) in the equation \(a=c\) or
\(c=a\); then infer \(E(c(d/y)/x)\).  Second shift: from \(b(c/x)\doteq d\)
and \(c\doteq a\) infer \(b(a/x)\doteq d\), substituting at \(x\) in
\(b\doteq d\); then infer \(E(d/y)\).  The side conditions that the
substitution variables do not occur in the substituted terms make these
literal instances.

Singularity.  Let \(\#_x(e)\) be the number of occurrences of \(x\) in
\(e\).  In a switch, the left shape is singular iff
\(\#_x(E)=\#_y(E)=1\) (the substituted terms contain neither variable), which
is also the condition for the right shape.  In the first shift, the upper
left inference needs \(\#_x(E)=1\) and the lower one
\(\#_y(E(c/x))=\#_x(E)\cdot\#_y(c)=1\); the right shape needs
\(\#_y(c)=1\) (as \(y\) does not occur in \(a\)) and \(\#_x(E)=1\).  Both
conditions say \(\#_x(E)=\#_y(c)=1\).  In the second shift, the left shape
needs \(\#_x(E(b/y))=\#_y(E)\cdot\#_x(b)=1\) and \(\#_y(E)=1\), and the right
shape needs \(\#_x(b)=1\) (as \(x\) does not occur in \(d\)) and
\(\#_y(E)=1\); again the conditions coincide.  Every other inference is
unchanged.  So a move and its inverse both preserve singularity, and
induction on the number of moves proves the claim in both directions.
\end{proof}
\end{annexitem}

\begin{annexitem}{annex:p02-w02-p02-028-st-nesting}{thm:p02-w02-p02-028-st-nesting}{Nested inferences}
\begin{proof}
A replacement determines its inference of form (1) up to the name of the
rule variable, which plays no role; so it suffices to exhibit the shapes with
the replaced and inserted terms in the right places.  For an equation \(e\)
and a position \(\pi\) of \(e\), write \(e|_\pi\) for the subterm at \(\pi\)
and \(e[c]_\pi\) for the result of putting \(c\) there.

(ii)  Let \(I\) replace \(a\) by \(c\) at \(\rho\) in its main premise \(e\),
with side premise \(a\doteq c\), and let \(J\) replace \(b\) by \(d\) at
\(\pi\) in \(e'=e[c]_\rho\), with side premise \(b\doteq d\).  Let \(T_1\),
\(T_2\), \(T_3\) be the subcomputations of \(e\), of \(a\doteq c\) and of
\(b\doteq d\).
\begin{itemize}
  \item \(\rho\) and \(\pi\) disjoint.  Then \(e|_\pi=e'|_\pi=b\).  Let
        \(E\) be \(e'\) with new variables \(x\) at \(\rho\) and \(y\) at
        \(\pi\).  Then \(e=E(a/x,b/y)\), \(e'=E(c/x,b/y)\), and the
        conclusion of \(J\) is \(E(c/x,d/y)\): the left shape of a switch.
  \item \(\pi=\rho\sigma\).  Then \(c|_\sigma=b\).  Let \(c_0\) be \(c\) with
        a new \(y\) at \(\sigma\) and \(E\) be \(e\) with a new \(x\) at
        \(\rho\).  Then \(e=E(a/x)\), \(I\)'s side premise is
        \(a\doteq c_0(b/y)\), \(e'=E(c_0(b/y)/x)\), and \(J\) concludes
        \(E(c_0(d/y)/x)\): the left shape of the first shift.  (This includes
        \(\sigma\) empty, with \(c_0=y\).)
  \item \(\rho=\pi\sigma\) with \(\sigma\) nonempty.  Then \(b=e'|_\pi\) has
        \(c\) at \(\sigma\).  Let \(b_0\) be \(b\) with a new \(x\) at
        \(\sigma\) and \(E\) be \(e'\) with a new \(y\) at \(\pi\).  As \(e\)
        and \(e'\) differ only at \(\rho\), \(e|_\pi=b_0(a/x)\), so
        \(e=E(b_0(a/x)/y)\), \(e'=E(b_0(c/x)/y)\), \(J\)'s side premise is
        \(b_0(c/x)\doteq d\), and \(J\) concludes \(E(d/y)\): the left shape
        of the second shift.
\end{itemize}
Two positions are either disjoint or comparable, so these cases are
exhaustive.

(iii)  Let \(J\) replace \(u\) by \(v\) at \(\pi\) in its main premise
\(e\), so that its side premise is \(u\doteq v\), and let this side premise
be the conclusion of \(I\), which replaces \(g\) by \(h\) at \(\rho\) in its
main premise \(f\), with side premise \(g\doteq h\).  Let \(E\) be \(e\) with
a new variable at \(\pi\).  The position \(\rho\) lies inside the side of the
equation \(u\doteq v\) that is \(v\), or inside the side that is \(u\), say
at the relative position \(\sigma\).
\begin{itemize}
  \item Inside \(v\).  Then \(v|_\sigma=h\); let \(c_0\) be \(v\) with a new
        \(y\) at \(\sigma\).  So \(f\) is \(u\doteq c_0(g/y)\) (with the
        orientation of \(u\doteq v\)), \(I\) infers \(u\doteq c_0(h/y)\)
        from it and \(g\doteq h\), and \(J\) infers \(E(c_0(h/y)/x)\) from
        \(E(u/x)\) and \(u\doteq c_0(h/y)\).  This is the right shape of the
        first shift, with \(a=u\), \(b=g\), \(d=h\) and \(c=c_0\).
  \item Inside \(u\).  Then \(u|_\sigma=h\); let \(b_0\) be \(u\) with a new
        \(x'\) at \(\sigma\).  So \(f\) is \(b_0(g/x')\doteq v\), \(I\)
        infers \(b_0(h/x')\doteq v\) from it and \(g\doteq h\), and \(J\)
        infers \(E(v/y)\) from \(E(b_0(h/x')/y)\) and \(b_0(h/x')\doteq v\).
        This is the right shape of the second shift, with \(b=b_0\),
        \(a=h\), \(c=g\) and \(d=v\); its premise \(c\doteq a\) is
        \(g\doteq h\) read in the other orientation, which the ambiguous
        notation \(\doteq\) allows.
\end{itemize}
\end{proof}
\end{annexitem}

\begin{annexitem}{annex:p02-w02-p02-028-st-normal-forms}{thm:p02-w02-p02-028-st-normal-forms}{The forms \(T_R\) and \(T_L\)}
\begin{proof}
\emph{(P1) Finitely many leftness values.}  By
Lemma~\ref{thm:p02-w02-p02-028-st-moves-basic}, moves and their inverses
preserve singularity and length, and they keep the leaves of
\(T_1,T_2,T_3\), so the number \(n\) of leaves is preserved.  Hence all
computations obtained from \(T\) by moves and inverse moves have leftness in
the finite set of \(n\)-tuples with entries between \(1\) and
\(\operatorname{lh}(T)\).

\emph{(P2) A shift decreases leftness.}  In the left shape of a shift, every
maximal path through a leaf of \(T_1\) passes through the root of the shape,
the conclusion of the upper inference, and then \(T_1\); in the right shape
\(T_1\) stands directly above the root.  So each such path is shorter by one,
paths through leaves outside the shape are unchanged, and the leaves of
\(T\) to the left of the first leaf of \(T_1\) are exactly the leaves to the
left of the shape, because \(T_1\) is the leftmost part of the shape.  Leftness therefore
decreases strictly, at the index of the first leaf of \(T_1\).  For the same
reason the inverse of a shift increases leftness strictly.

\emph{(P3) Switches on a spine.}  Let \(J_1,\ldots,J_m\) be a spine (bottom
first) and \(T_0\) the subcomputation of its top equation.  A switch of
\(J_{i+1}\) over \(J_i\) yields two inferences nested to the left that carry
the replacements of \(J_{i+1}\) and \(J_i\) in the other order, at the same
positions with the same terms, while \(T_1\), the subcomputation of the main
premise of \(J_{i+1}\), is unchanged.  So the spine keeps its length, its
top subcomputation \(T_0\) and the depth of every leaf of \(T_0\); the side
subcomputations are permuted without being changed, and no other spine is
altered.  A switch is its own inverse: its right shape is the left shape of a
switch, which returns the original.

\emph{Part 1.}  Repeat the following while some spine of the current
computation has two inferences at comparable positions.  Choose such a pair
\(J_k\) (below) and \(J_j\) (above) with \(j-k\) least.  Every \(J_i\) with
\(k<i<j\) has its position disjoint from both, by minimality, and a
replacement at a position disjoint from another leaves the other's position
and subterm in place.  By Lemma~\ref{thm:p02-w02-p02-028-st-nesting}(ii),
\(J_j\) and \(J_{j-1}\) have the left shape of a switch; switching them moves
the replacement of \(J_j\) one place down the spine.  After \(j-k-1\) such
switches the replacements of \(J_j\) and \(J_k\) are adjacent and nested to
the left at comparable positions, so by the same lemma they have the left
shape of a shift, which is applied.  All these moves involve inferences of
the one spine, that is, of the leftmost path of the subcomputation ending in
its bottom inference.  By (P3) the switches keep the depth of the leaves of
\(T_0\), and every leaf left of the first leaf of \(T_0\) lies outside that
subcomputation and keeps its path; by (P2) the shift then shortens the path of
the first leaf of \(T_0\), which is the first leaf of its \(T_1\).  So each
repetition decreases leftness, and by (P1) the repetitions stop.  The result
is \(T_R\); \(T\) reduces to it, and in \(T_R\) the positions on each spine
are pairwise disjoint.

Call this property (Q).  A shift needs, in its left shape, two inferences
nested to the left at comparable positions: in the first shift the lower
inference replaces inside the term just inserted by the upper one, in the
second the lower inference replaces a term containing it.  So no shift
applies to a computation with (Q).  A switch preserves (Q), by (P3).  Hence,
by induction along any sequence of moves beginning with \(T_R\), every
computation in it has (Q) and every move in it is a switch.  This is
[2.11(a)].

\emph{Part 2.}  Repeat the following while some inference \(J\) of form (1)
has as its side premise the conclusion of an inference \(I\) of form (1).  By
Lemma~\ref{thm:p02-w02-p02-028-st-nesting}(iii), the subcomputation ending
in \(J\) has the right shape of a shift; replace it by the left shape.  By
(P2) leftness increases strictly, so by (P1) the repetitions stop.  The
result is \(T_L\), and reading the repetitions backwards, \(T_L\) reduces to
\(T\) by shifts.  Call (Q\('\)) the property that no side premise of an
inference of form (1) is the conclusion of an inference of form (1); \(T_L\)
has it.  A switch preserves (Q\('\)): the side premises of its two new
inferences are the roots of \(T_3\) and \(T_2\), which were side premises
before, and the new middle equation is a main premise.  As a switch is its
own inverse, a switch also preserves the failure of (Q\('\)).  The right
shape of a shift fails (Q\('\)).  Now let a sequence of moves end in \(T_L\).
If it contained a shift, the computation just after the last shift would
fail (Q\('\)), and so, by the switches that follow, would \(T_L\); this is
impossible.  So the sequence consists of switches.  This is [2.11(b)].

\emph{Part 3.}  Let a spine of \(T_R\) carry, from the top, replacements of
\(a_1\) by \(b_1\), \ldots, \(a_n\) by \(b_n\) at pairwise disjoint positions
\(\pi_1,\ldots,\pi_n\).  A replacement at one position leaves the subterms
at disjoint positions in place, so the top equation has \(a_i\) at \(\pi_i\)
for every \(i\).  Let \(E\) be the top equation with new distinct variables
\(x_i\) at \(\pi_i\).  Then the top equation is
\(E(a_1/x_1,\ldots,a_n/x_n)\), the conclusion of the \(i\)-th inference is
\(E(b_1/x_1,\ldots,b_i/x_i,a_{i+1}/x_{i+1},\ldots,a_n/x_n)\), and its side
premise is \(a_i\doteq b_i\): the left form of
Figure~\ref{fig:p02-w02-p02-028-st-fragments}.  In \(T_L\), by (Q\('\)) two
inferences of an \(H\) fragment can be connected only by nesting to the
left, so an \(H\) fragment is a chain of inferences nested to the left, that
is, one spine, and its side premises are not conclusions of inferences of
form (1): the right form of the figure.  If \(T\) is an \(H\) computation, so
is \(T_L\); every equation occurrence of \(T_L\) that is not a leaf is the
conclusion of an inference of form (1), so its side premises are leaves,
which are axioms or assumptions, and the spine through the root contains
every inference.
\end{proof}
\end{annexitem}

\begin{annexitem}{annex:p02-w02-p02-028-st-prop-3-1}{thm:p02-w02-p02-028-st-prop-3-1}{Direct computations from simple sets}
\begin{proof}
Run Method~\ref{def:p02-w02-p02-028-st-closure}.  Write \(a_n,b_n\) for the
pair chosen at stage \(n+1\), so \(S_{n+1}=S_n(a_n/b_n)\).  In both cases
\(a_n\) and \(b_n\) are left sides of simple equations, hence variables;
\(S_{n+1}\) is simple and has one variable fewer than \(S_n\), so the
procedure stops at some \(n_0\leq k\) (if \(k=0\), \(n_0=1\)).  Let
\(\sigma_n\) be the composite \((a_1/b_1)\cdots(a_{n-1}/b_{n-1})\) and
\(\sigma=\sigma_{n_0}\); then \(S_n=S_1\sigma_n\), \(x_n=x\sigma_n\) and
\(y_n=y\sigma_n\), and \(\sigma\) sends variables to variables.

\emph{The auxiliary computations.}  \(C_1\) consists of the one-node
computations of the members of \(S_1\) (assumptions, and axioms \(x=x\)).
Given \(C_n\), with a computation in \(C_n\) of each member of \(S_n\):
in Case~1, from \(b_n=c\) and \(a_n=c\) rule (1), with the equation
\(b_n=w\) substituted at the fresh \(w\), gives \(b_n=a_n\); in Case~2, from
\(c=sa_n\) and \(c=sb_n\) rule (1), with \(w=sa_n\) at \(w\), gives
\(sb_n=sa_n\), and rule (2) gives \(b_n=a_n\).  This is \(T^*_n\).  For each
\(E\in S_n\) with computation \(T\in C_n\), rule (1) applied to \(E\), read as
\(E(b_n/w)\) with \(w\) fresh at the occurrences of \(b_n\), and to
\(b_n\doteq a_n\) from \(T^*_n\) gives \(E(a_n/b_n)\); these form \(C_{n+1}\),
and every member of \(S_{n+1}\) is such an \(E(a_n/b_n)\).  Let \(h_n\) be
the largest number of equation occurrences on a branch of a member of
\(C_n\).  Then \(h_1=1\), \(T^*_n\) has branches of at most \(h_n+2\)
occurrences, and \(h_{n+1}\leq h_n+3\); so \(h_n\leq3n-2\) and the branches
of \(T^*_n\) have at most \(3n\) occurrences.  All computations so far use
only rules (1) and (2), from \(S\).

\emph{Directness of the material.}  Every term occurring in these
computations occurs in some \(S_n\).  If \(t\) occurs in \(S_{n+1}\), then
\(t=t'(a_n/b_n)\) for some \(t'\) in \(S_n\); writing \(t'=t_0(b_n/w)\) with
\(w\) fresh, the first measure condition and \(M(a_n)\leq M(b_n)\) give
\(M(t)=M(t_0(a_n/w))\leq M(t_0(b_n/w))=M(t')\).  By induction, every term
occurring in any \(S_n\) has measure at most that of a term of \(S_1\), whose
terms occur in \(S\).  The same argument gives \(M(x_n)\leq M(x)\) and
\(M(y_n)\leq M(y)\).

\emph{Properties of \(S_{n_0}\) that need no hypothesis.}  (i) A term
\(fz_1\cdots z_m\) occurs in a simple set only as a right side, so some
\(z=fz_1\cdots z_m\) belongs to \(S_{n_0}\); two such with distinct left
sides would trigger Case~1.  (ii) and (iii) hold likewise for \(sz_1\) and
\(0\).  (iv) If \(z_1=z_2\in S_{n_0}\) with \(z_1\) not \(z_2\), then together
with \(z_2=z_2\in S_{n_0}\) (the image of an axiom \(x=x\)) Case~1 applies.
(v) If \(z=sz_1\) and \(z=sz_2\) belong to \(S_{n_0}\) with \(z_1\) not \(z_2\),
Case~2 applies.

\emph{The read-off interpretation.}  Suppose \(S_{n_0}\) also satisfies (vi)
and (vii).  Let \(V\) be its set of variables and \(N=\abs V+1\).  By (v),
\(z\mapsto z'\) for \(z=sz'\in S_{n_0}\) is a partial function on \(V\); by
(vii) it has no cycle; by (ii) no variable is the image of two variables.
So \(V\) is a disjoint union of chains \(r,z_1,z_2,\ldots\) with
\(z_1=sr\), \(z_2=sz_1\), and so on, each starting at a variable \(r\) that is
not a left side of an \(s\)-equation.  By (iii) at most one \(r_0\) has
\(r_0=0\in S_{n_0}\), and by (vi) it starts a chain.  Number the chain starts
\(r_0,r_1,\ldots\) (with \(r_0\) first when it exists) and set
\(I_0(z)=Nj+\text{(position of \(z\) in the chain of \(r_j\))}\).  Positions
are below \(N\), so \(I_0\) is injective, \(I_0(r_0)=0\), and
\(I_0(z)=I_0(z')+1\) when \(z=sz'\in S_{n_0}\).  For a function variable
\(f\), let \(I_0(f)(I_0(z_1),\ldots,I_0(z_m))=I_0(z)\) when
\(z=fz_1\cdots z_m\in S_{n_0}\), well defined by injectivity and (i), and
\(0\) at all other arguments.  Each equation of \(S_{n_0}\) is true under
\(I_0\): \(z=0\) as \(z\) is \(r_0\); \(z_1=z_2\) by (iv); \(z=sz'\) and
\(z=fz_1\cdots z_m\) by construction.  Put \(I(v)=I_0(v\sigma)\) for
variables \(v\) of \(S\) and \(I(f)=I_0(f)\).  Then \(I(t)=I_0(t\sigma)\) for
terms of \(S\), and \(e\sigma\in S_{n_0}\) for \(e\in S\), so \(I\) satisfies
\(S\).

\emph{(i) The analytic case.}  If \(x\) is \(y\), the axiom \(x=x\) is the
required computation.  Otherwise \(x\) and \(y\) occur in \(S\): if, say,
\(x\) did not, a model of \(S\) could be changed at \(x\) to make
\(x\neq y\).  Let \(I'\) satisfy \(S\).  By (T1) every \(T^*_n\) is true under
\(I'\), so \(I'(a_n)=I'(b_n)\), \(I'(v\sigma)=I'(v)\) for all variables, and
\(I'\) satisfies \(S_{n_0}\).  Hence (vi) holds, since \(0\neq I'(z_2)+1\), and
(vii) holds, since \(I'(z_1)=I'(z_1)+m\) is impossible.  The read-off \(I\)
satisfies \(S\), and validity gives \(I_0(x\sigma)=I(x)=I(y)=I_0(y\sigma)\);
by injectivity \(x_{n_0}=x\sigma\) is \(y_{n_0}=y\sigma\).  Now build
\(U_{n_0},\ldots,U_1\): \(U_{n_0}\) is the axiom \(x_{n_0}=y_{n_0}\).  For
\(i<n_0\), if neither \(x_i\) nor \(y_i\) is \(b_i\), then
\(x_{i+1}=y_{i+1}\) is \(x_i=y_i\) and \(U_i=U_{i+1}\).  Otherwise let \(e\) be
\(x_i=y_i\) with each side that is \(b_i\) replaced by a fresh \(w\); then
\(e(a_i/w)\) is \(x_{i+1}=y_{i+1}\), and rule (1) applied to \(U_{i+1}\) and
\(T^*_i\) gives \(e(b_i/w)\), that is \(x_i=y_i\); this is \(U_i\).  \(U_1\)
computes \(x=y\) from \(S\).  By (T2) its length is at most
\(1+\sum_{i=1}^{n_0-1}2^{3i}=1+(8^{n_0}-8)/7\leq8^{n_0}\leq2^{3k}\).  Its terms
are \(x_i,y_i\) and terms of the \(S_n\), so it is \(M\)-direct.

\emph{(ii) The unsatisfiable case.}  If (vi) and (vii) both held, the
read-off \(I\) would satisfy \(S\).  So one fails.  If (vi) fails, with
\(z_1=sz_2\) and \(z_1=0\) in \(S_{n_0}\), rule (1) applied to their
computations in \(C_{n_0}\), substituting at \(w\) in \(w=sz_2\), gives
\(0=sz_2\); \(z_2\) occurs in \(S\), the terms are \(M\)-bounded as before
(\(sz_2\) occurs in the conclusion), and the length is at most
\(1+2(2^{3k-2}-1)<2^{4k}\).  If (vii) fails, (v) makes \(z\mapsto z'\) a
partial function, so there is a cycle \(z_1=sz_m,z_2=sz_1,\ldots,z_m=sz_{m-1}\)
in \(S_{n_0}\) with distinct \(z_j\) and \(m\geq1\); indices are read modulo
\(m\).  Its \(m\) equations are images under \(\sigma\) of distinct
equations of the form \(x_1=sx_2\) of \(S\), so \(m\) is at most their number,
hence \(m\leq\operatorname{card}(S)\); also \(m\leq k\).  Choose \(i\) with
\(M(z_i)\) maximal among \(z_1,\ldots,z_m\).  Let \(V_1\) be the computation
in \(C_{n_0}\) of \(z_i=sz_{i-1}\), and obtain \(V_{q+1}\) from \(V_q\), which
computes \(z_i=s^qz_{i-q}\), by rule (1) at \(w\) in \(z_i=s^qw\) with the
computation of \(z_{i-q}=sz_{i-q-1}\).  \(V_m\) computes \(z_i=s^mz_i\).
Every term \(s^qz_j\) (\(q\leq m\)) in it satisfies
\(M(s^qz_j)\leq M(s^qz_i)\leq M(s^mz_i)\), by the first measure condition
(as \(M(z_j)\leq M(z_i)\)) and the second (with \(c=s^{m-q}w\)); and
\(s^mz_i\) occurs in the conclusion.  The printed proof writes this relation
with an equality sign \cite[p.~906]{Statman1977}; the inequality is what the
measure conditions give and what directness requires.  The length is at most
\(m\,2^{3k-2}+m-1\leq k\,2^{3k}\leq2^{4k}\), as \(k\leq2^k\).
\end{proof}
\end{annexitem}

\begin{annexitem}{annex:p02-w02-p02-028-st-prop-3-2}{thm:p02-w02-p02-028-st-prop-3-2}{Direct computations from arbitrary sets}
\begin{proof}
(i)  For each occurrence \(t\) of a term in \(\bigwedge S\to a=b\) take a new
variable \(x_t\).  Let \(S_1'\) contain \(x_t=0\) if \(t\) is \(0\),
\(x_t=y\) if \(t\) is the variable \(y\), \(x_t=sx_{t'}\) if \(t\) is
\(st'\), and \(x_t=fx_{t_1}\cdots x_{t_m}\) if \(t\) is \(ft_1\cdots t_m\)
(the \(x_{t'}\) and \(x_{t_i}\) for the immediate sub-occurrences); let
\(S_2'=\{x_{t_1}=x_{t_2}:t_1=t_2\in S\}\) for the top occurrences.  The set
\(S'=S_1'\cup S_2'\) is simple.  Let \(\rho\) send each \(x_t\) to the term
at \(t\) and fix every other variable.

\emph{(\(\alpha\)) Any interpretation satisfying \(S_1'\) gives \(x_t\) the
value of \(t\)}, by induction on \(t\).  Hence a model of \(S'\) is a model of
\(S\), and then of \(a=b\), so of \(x_a=x_b\); and a model of \(S\) extended
by \(x_t\mapsto\) value of \(t\) satisfies \(S'\).  So
\(\bigwedge S'\to x_a=x_b\) is analytic.  \(M^*(c)=M(c\rho)\) is a measure by
Lemma~\ref{thm:p02-w02-p02-028-st-measure-transfer} (\(\rho\) replaces only
individual variables).  Theorem~\ref{thm:p02-w02-p02-028-st-prop-3-1}(i)
gives an \(M^*\)-direct computation \(T'\) of \(x_a=x_b\) from \(S'\), by
rules (1) and (2), of length at most \(2^{3k}\), where \(k\) counts the
variables of \(S'\): one per term occurrence, each of which begins at a
distinct symbol occurrence, plus the original variables; so
\(k\leq2(\operatorname{lh}(S)+\operatorname{lh}(a=b))\) and
\(2^{3k}\leq2^{9(\operatorname{lh}(S)+\operatorname{lh}(a=b))}\).

Apply \(\rho\) to every equation of \(T'\), after renaming the substitution
variable of each rule-(1) inference to one occurring nowhere else.
Assumptions in \(S_1'\) become axioms \(0=0\), \(y=y\), \(st'=st'\),
\(ft_1\cdots t_m=ft_1\cdots t_m\); assumptions in \(S_2'\) become members of
\(S\); axioms stay axioms.  An inference
\(E(c/w),\,c\doteq d\,/\,E(d/w)\) becomes
\((E\rho)(c\rho/w),\,c\rho\doteq d\rho\,/\,(E\rho)(d\rho/w)\), again an
instance of rule (1), and rule (2) is preserved.  So \(\rho T'\) computes
\(a=b\) from \(S\) with the same length.  Each term of \(\rho T'\) is
\(e\rho\) for a term \(e\) of \(T'\), and some \(e'\) in \(x_a=x_b\) or \(S'\)
has \(M(e\rho)\leq M(e'\rho)\); \(e'\rho\) is a term occurring in \(a=b\) or
\(S\), or a subterm \(t\) of such a term \(t_0\), and then \(M(t)\leq M(t_0)\)
by the second measure condition applied to \(t_0=c(t/w)\) with \(c\) the
context of that occurrence.  So \(\rho T'\) is \(M\)-direct.

(ii)  Flatten \(\bigwedge S\) alone.  If \(S'\) had a model, (\(\alpha\))
would give one of \(S\); so \(S'\) is unsatisfiable, and
Theorem~\ref{thm:p02-w02-p02-028-st-prop-3-1}(ii) gives an \(M^*\)-direct
computation of \(x=s^px\) or \(0=sx\) from \(S'\) of length at most
\(2^{4k}\leq2^{8\operatorname{lh}(S)}\leq2^{12\operatorname{lh}(S)}\), with
\(p\) at most the number of equations \(x_{st'}=sx_{t'}\), hence at most the
number of occurrences of \(s\) in \(S\), so \(p\leq\operatorname{lh}(S)\).
Applying \(\rho\) as in (i) gives an \(M\)-direct computation of
\(a=s^pa\) or \(0=sa\) from \(S\) with \(a=x\rho\) occurring in \(S\).
\end{proof}
\end{annexitem}

\begin{annexitem}{annex:p02-w02-p02-028-st-completeness}{thm:p02-w02-p02-028-st-completeness}{Completeness of \(K\), second proof}
\begin{proof}
(1) Write \(E\) as \(a=b\).  The non-logical symbols of
\(\bigwedge S\to E\) are exactly the symbols of its terms, so
\(p=\operatorname{lh}(S)+\operatorname{lh}(a=b)\), and
Theorem~\ref{thm:p02-w02-p02-028-st-prop-3-2}(i) with \(M=\operatorname{lh}\)
gives the computation.

(2) If \(\bigwedge S\) is satisfiable, item~1 applies.  Otherwise
Theorem~\ref{thm:p02-w02-p02-028-st-prop-3-2}(ii) gives a computation of
\(c=s^qc\) with \(q\geq1\), or of \(0=sc\); one application of rule (\(4_q\))
or (3) concludes \(E\).

(3) Axioms and rules (1), (2) preserve truth under every interpretation by
(T1).  The premise \(0=sa\) of rule (3) is false under every interpretation,
since \(0\) is not a successor, and the premise \(a=s^na\) (\(n\geq1\)) of
rule (\(4_n\)) is false, since \(m\neq m+n\).  So, by induction on
computations, every equation computed from \(S\) is true under every
interpretation satisfying \(S\); with item~2 this gives the equivalence.

(4) In the flattened set \(S'\) of the previous proof, there are at most
\(p\) equations from \(S_1'\) and at most \(\operatorname{card}(S)\leq p\)
from \(S_2'\), and \(k\leq2p\) variables, so each stage \(S_n\) of the
identification procedure has at most \(4p\) members.  List the members of
\(S_1\); then, for \(n=1,\ldots,n_0-1\), the at most two equations of
\(T^*_n\) not yet listed, followed by the members of \(S_{n+1}\), each
obtained by rule (1) from a listed member of \(S_n\) and \(b_n=a_n\); then the
equations \(x_i=y_i\) for \(i=n_0,\ldots,1\).  Every equation is an axiom,
an assumption, or the conclusion of a rule whose premises occur earlier, and
the list contains every equation of the tree \(U_1\).  Its length is at most
\(n_0\cdot4p+2n_0+n_0\leq8p^2+6p\leq16p^2\) for \(p\geq1\).  Applying
\(\rho\) to each member preserves these properties and the length.
\end{proof}
\end{annexitem}

\begin{annexitem}{annex:p02-w02-p02-028-st-decision}{thm:p02-w02-p02-028-st-decision}{Polynomial-time decision}
\begin{proof}
\emph{Algorithm.}  If \(a\) is \(b\), answer ``valid''.  Otherwise form the
flattened set \(S'\) of \(\bigwedge S\to a=b\) and its target \(x_a=x_b\),
and run Method~\ref{def:p02-w02-p02-028-st-closure} with
\(M=\operatorname{lh}\), breaking ties by a fixed order of variables.  If the
final stage violates (vi) or (vii), answer ``valid''.  Otherwise answer
``valid'' iff \(x_a\sigma\) is \(x_b\sigma\).

\emph{Correctness.}  By (\(\alpha\)) in the previous proof, \(S'\) is
satisfiable iff \(S\) is.  The proof of
Theorem~\ref{thm:p02-w02-p02-028-st-prop-3-1} shows that, whatever measure
guides the choices, \(S'\) is unsatisfiable iff the final stage violates (vi)
or (vii): a model of \(S'\) forces both, and if both hold the read-off
interpretation is a model.  If \(S\) is unsatisfiable, the implication is
valid.  If \(S\) is satisfiable and the implication is valid, then
\(\bigwedge S'\to x_a=x_b\) is analytic and the same proof shows that
\(x_a\sigma\) is \(x_b\sigma\).  Conversely, if \(x_a\sigma\) is
\(x_b\sigma\), the computation \(U_1\) of that proof and the map \(\rho\) give
a computation of \(a=b\) from \(S\), so the implication is valid by
Theorem~\ref{thm:p02-w02-p02-028-st-completeness}(3).

\emph{Cost.}  Let \(p=\operatorname{lh}(S)+\operatorname{lh}(a=b)\).  \(S'\)
is produced in one pass and has at most \(2p\) equations; \(S_1\) has at most
\(4p\), each of length at most \(p+1\) variable names.  There are at most
\(k\leq2p\) stages.  A stage compares the right sides of all pairs of
equations, at most \((4p)^2\) comparisons of at most \(p+1\) names each, and
then rewrites one name throughout, at most \(4p(p+1)\) name operations.  The
final check of (vi) compares pairs, and the check of (vii) follows the
partial map \(z\mapsto z'\) at most \(k\) steps from each variable.  In total
at most \(O(p^4)\) operations on variable names, each of
\(O(\log p)\) bits.  A deterministic Turing machine working on a code of the
formula, whose length is at least \(p\), performs each such operation in time
polynomial in the code length, so the whole decision runs in polynomial time.
\end{proof}
\end{annexitem}

\begin{annexitem}{annex:p02-w02-p02-028-st-J-route}{thm:p02-w02-p02-028-st-J-route}{The model-theoretic route}
\begin{proof}
(1) (1) holds as equality is identity on \(\mathbb N\).  For (2) and (3), if
\(x\) and \(y\) denote the same number, \(f\) and \(g\) take the same values
there.  (4) holds as successor is injective.  The antecedents of (5) and
(\(6_n\)) are never true, as \(0\) is no successor and \(m\neq m+n\) for
\(n\geq1\).

(2) Suppose \(\bigwedge S\to a=b\) is not valid.  Then \(S\) is satisfiable
and, by the correctness part of the previous proof, the final stage of the
procedure on the flattened set satisfies (vi) and (vii) and
\(x_a\sigma\) is not \(x_b\sigma\).  The read-off interpretation \(I\) of the
proof of Theorem~\ref{thm:p02-w02-p02-028-st-prop-3-1} satisfies \(S'\), and,
\(I_0\) being injective, \(I(x_a)\neq I(x_b)\); by (\(\alpha\)) it satisfies
\(S\) and falsifies \(a=b\).  Its function variables have value \(0\) at all
but finitely many arguments, and its individual values are below
\(N\cdot N\leq(k+1)^2\), \(N=\abs V+1\leq k+1\).  So if the implication is not
valid, it fails under an interpretation in which every individual value,
every argument tuple at which a function variable is nonzero, and every
value is below \((k+1)^2\).  There are finitely many such interpretations of
the finitely many symbols of the implication, and each is checked by
evaluation, so validity is decidable.

(3) Let \(\mathfrak A\) be a general model of \(J\) as described, with an
assignment under which every member of \(S\) holds and \(E\) fails.  Suppose
\(\bigwedge S\to E\) were valid.  By
Theorem~\ref{thm:p02-w02-p02-028-st-completeness}(2), \(E\) has a
computation from \(S\) with rules (1)--(\(4_n\)).  Every step preserves truth
in \(\mathfrak A\): an axiom \(a=a\) by axiom (1).  For rule (1) with
\(E'=(l=r)\), let \(f=\lambda x\,l\) and \(g=\lambda x\,r\); the premise
\(E'(a/x)\) says \(fa=ga\), and axiom (2), for \(a=b\), or axiom (3), for
\(b=a\), gives \(fb=gb\), which is \(E'(b/x)\).  Rule (2) is preserved by
axiom (4), rule (3) by axiom (5) and rule (\(4_n\)) by axiom (\(6_n\)),
instantiating \(y,z\) by the values of the conclusion's sides.  So \(E\)
would hold in \(\mathfrak A\), a contradiction.  Hence the implication is not
valid, and item~2 provides a countermodel of the stated form.
\end{proof}
\end{annexitem}

\begin{annexitem}{annex:p02-w02-p02-028-st-prop-3-3-2}{thm:p02-w02-p02-028-st-prop-3-3-2}{From the subterm property to \(M\)-directness}
\begin{proof}
Here a term \emph{occurs in} \(E\) or \(S\) when it is a subterm of a term of
\(E\) or of a member of \(S\).  Moves keep the leaves of a computation and
every inference of form (2) with its premise and conclusion, because they
rearrange only two inferences of form (1) above the unchanged
subcomputations \(T_1,T_2,T_3\) and below an unchanged root.  So every
computation to which \(T\) reduces is singular
(Lemma~\ref{thm:p02-w02-p02-028-st-moves-basic}), computes \(E\) from \(S\)
with the same length, and has the subterm property.

\emph{The switching process.}  Let \(T_R\) be as in
Theorem~\ref{thm:p02-w02-p02-028-st-normal-forms}(1).  Call an inference of
form (1) that replaces \(a\) by \(b\) \emph{ascending} if \(M(a)<M(b)\), and
\emph{non-ascending} if \(M(b)\leq M(a)\).  An \emph{inversion} is a pair of
inferences on one spine with the ascending one above the non-ascending one.
While some spine has an ascending inference immediately above a
non-ascending one, switch them.  The switch applies: positions on each spine
of \(T_R\) are pairwise disjoint, switches preserve this
(Theorem~\ref{thm:p02-w02-p02-028-st-normal-forms}, proof of part~1), and two
adjacent inferences at disjoint positions have the left shape of a switch
(Lemma~\ref{thm:p02-w02-p02-028-st-nesting}(ii)).  The switch exchanges the
two replacements, each keeping its terms and so its kind, and changes no
other spine (step (P3) of that proof).  So the number of inversions decreases
by exactly one, and the process stops after at most as many switches as
\(T_R\) has inversions.  Call the result \(T_M\).  On each spine of \(T_M\)
no ascending inference stands immediately above a non-ascending one, so all
non-ascending inferences lie above all ascending ones.  An ascending
replacement of \(a\) by \(b\) at \(x\) immediately followed by a
non-ascending replacement of \(c\) by \(d\) at a disjoint \(y\) is exactly
the configuration displayed in [3.3.2], with \(M(a)<M(b)\) and
\(M(d)\leq M(c)\); so \(T_M\) contains none.  \(T\) reduces to \(T_R\) and
\(T_R\) to \(T_M\).

\emph{\(M\)-directness.}  Call a term \emph{bounded} if its measure is at most
that of some term occurring in \(E\) or \(S\), and an equation bounded if
both sides are.
\begin{itemize}
  \item[(B0)] A subterm \(t\) of a bounded term \(t'\) is bounded: write
        \(t'=t_0(t/w)\) with \(w\) a new variable occurring in \(t_0\); the
        second measure condition gives \(M(t)\leq M(t')\).  Every term
        occurring in \(T_M\) is a subterm of a side of one of its equations,
        so it suffices to show that every equation occurrence of \(T_M\) is
        bounded.
  \item[(B1)] If an inference replaces \(a\) by \(b\) inside the side
        \(l=l_0(a/w)\) of its main premise, the new side is \(l_0(b/w)\), the
        other side is unchanged, and by the first measure condition
        \(M(l_0(b/w))\leq M(l_0(a/w))\) when the inference is non-ascending
        and \(M(l_0(a/w))\leq M(l_0(b/w))\) when it is ascending.
  \item[(B2)] Let a spine have equations \(e_0\) (its top equation),
        \(e_1,\ldots,e_m\), where \(e_i\) is the conclusion of its \(i\)-th
        inference from the top, and let the first \(r\) inferences be the
        non-ascending ones.  By (B1), for \(i\leq r\) each side of \(e_i\)
        has measure at most that of the same side of \(e_0\), and for
        \(i\geq r\) at most that of the same side of \(e_m\).  So if \(e_0\)
        and \(e_m\) are bounded, every equation of the spine is.
  \item[(B3)] Every top equation is bounded.  It is not the conclusion of an
        inference of form (1), so it is an assumption, whose sides occur in
        \(S\); an axiom \(a=a\), where \(a\) occurs in \(E\) or \(S\) by the
        subterm property; or the conclusion \(a=b\) of an inference of form
        (2) from \(sa=sb\), where \(sa\) and \(sb\) occur in \(S\) by the weak
        subterm property, and hence so do \(a\) and \(b\).
  \item[(B4)] Every bottom equation is bounded, by induction on the number of
        equation occurrences on the path from the root of \(T_M\) to it.  The
        bottom equation of a spine is not a main premise of an inference of
        form (1).  If it is the root, it is \(E\).  If it is the premise
        \(sa=sb\) of an inference of form (2), its sides occur in \(S\).
        Otherwise it is the side premise \(u\doteq v\) of an inference \(J\)
        of form (1) on another spine, which replaces \(u\) by \(v\); \(u\) is a
        subterm of the main premise of \(J\) and \(v\) of its conclusion.  The
        bottom equation of that spine lies on the path from the root to
        \(J\)'s conclusion, hence is nearer the root, and is bounded by the
        induction hypothesis; its top equation is bounded by (B3); so all
        its equations are bounded by (B2), and \(u\), \(v\) are bounded by
        (B0).
\end{itemize}
Every equation occurrence of \(T_M\) is a leaf, the conclusion or premise of
an inference of form (2), or an equation of a spine.  Leaves and conclusions
of form (2) were treated in (B3), premises of form (2) in (B4), and spine
equations by (B2)--(B4).  So \(T_M\) is \(M\)-direct.
\end{proof}
\end{annexitem}

\begin{annexitem}{annex:p02-w02-p02-028-st-prop-3-4}{thm:p02-w02-p02-028-st-prop-3-4}{From the weak to the full subterm property}
\begin{proof}
As in the previous proof, moves keep the inferences of form (2), so every
computation obtained from \(T\) by moves or inverse moves is singular, has
the weak subterm property, and computes \(E\) from \(S\) with length
\(\operatorname{lh}(T)\).

\emph{Step 1.}  Take \(T_L\)
(Theorem~\ref{thm:p02-w02-p02-028-st-normal-forms}(2)).  Each of its \(H\)
fragments is one spine whose side premises are not conclusions of
inferences of form (1).

\emph{Step 2: left before right.}  Call an inference of form (1) a
\emph{left} or a \emph{right} replacement according as its position begins
with \(1\) or with \(2\).  While some spine has a right replacement
immediately above a left one, switch them; their positions are disjoint, so
Lemma~\ref{thm:p02-w02-p02-028-st-nesting}(ii) gives the left shape of a
switch.  Each switch keeps the kind of every replacement and the property
(Q\('\)) of the proof of
Theorem~\ref{thm:p02-w02-p02-028-st-normal-forms}, and lowers by one the
number of pairs on a spine with a right replacement above a left one; so the
process stops.  Then no \(H\) fragment has an inference from
\(a_1=b_1(c_1/x_1)\) and \(c_1\doteq d_1\) to \(a_1=b_1(d_1/x_1)\) above one
from \(a_2(c_2/x_2)=b_2\) and \(c_2\doteq d_2\) to \(a_2(d_2/x_2)=b_2\), which
is the condition of the chapter.  Each \(H\) fragment is now a spine with top
equation \(a_p=b_q\), whose first \(p-1\) inferences, with side premises
\(E_1,\ldots,E_{p-1}\), lead through \(a_{p-1}=b_q,\ldots\) to \(a_1=b_q\), and
whose last \(q-1\) inferences, with side premises
\(E_p,\ldots,E_{p+q-2}\), lead through \(a_1=b_{q-1},\ldots\) to \(a_1=b_1\)
(\(p,q\geq1\)).

\emph{Step 3: tops and bottoms.}  The bottom \(a_1=b_1\) of a fragment is
not a main premise of an inference of form (1) (the spine is maximal) nor a
side premise of one (by (Q\('\))).  So it is \(E\), at the root, or the
premise \(sa=sb\) of an inference of form (2), whose sides occur in \(S\).
In both cases \(a_1\) occurs in \(E\) or \(S\).  The top \(a_p=b_q\) is an
assumption, so \(a_p\) and \(b_q\) occur in \(S\); the conclusion of an
inference of form (2) from \(sa_p=sb_q\), so again they occur in \(S\); or an
axiom, so \(a_p\) is \(b_q\).

\emph{Step 4: replacing a fragment.}  Let the top be an axiom \(a_p=a_p\)
with \(a_p\) occurring in neither \(E\) nor \(S\).  For \(1\leq i<p\), the
inference with side premise \(E_{p-i}\) passes from \(a_{i+1}=b_q\) to
\(a_i=b_q\) by replacing \(u_i\) by \(v_i\) at a position \(1\sigma_i\); so
\(a_{i+1}=C_i(u_i/w)\) and \(a_i=C_i(v_i/w)\) for a term \(C_i\) with one
occurrence of a new variable \(w\), and \(E_{p-i}\) is \(u_i\doteq v_i\).
Replace the fragment by the spine that starts from the axiom \(a_1=a_1\) and,
for \(i=1,\ldots,p-1\), passes from \(a_1=a_i\) to \(a_1=a_{i+1}\) using
\(E_{p-i}\), read as \(v_i\doteq u_i\), by replacing \(v_i\) by \(u_i\) at
\(2\sigma_i\); this is the inference of form (1) with \(E\) equal to
\(a_1=C_i\), in which \(w\) occurs once.  It reaches \(a_1=a_p\), which is
\(a_1=b_q\) since \(a_p\) is \(b_q\), and continues with the last \(q-1\)
inferences of the old fragment unchanged.  The new spine has the same
bottom \(a_1=b_1\), uses each side premise \(E_1,\ldots,E_{p+q-2}\) once with
its subcomputation, and has \(p+q-1\) equations on the spine, as the old one
did; so the length is unchanged, the computation is singular, and its
inferences of form (2) are unchanged.  Its new axiom \(a_1=a_1\) has \(a_1\)
occurring in \(E\) or \(S\) by Step~3.  Fragments are disjoint and each
replacement changes only the spine equations of its own fragment, so all
such fragments can be replaced.

\emph{Step 5: deleting trivial inferences.}  An inference of form (1) whose
side premise is an axiom \(b=b\) replaces \(b\) by \(b\), so its conclusion
\(a=c\) is its main premise \(a=c\).  Replace the subcomputation ending in it
by the subcomputation of its main premise; this removes two equation
occurrences and changes no inference of form (2).  Repeat until no such
inference is left, and call the result \(T_S\).

\emph{Verification.}  \(T_S\) is a singular computation of \(E\) from
\(S\) with \(\operatorname{lh}(T_S)\leq\operatorname{lh}(T)\), and it has the
weak subterm property because its inferences of form (2) are those of \(T\).
Every leaf of \(T_S\) is the root, a premise of an inference of form (2), a
side premise of an inference of form (1), or the top equation of a spine.
Let it be an axiom \(a=a\).  As the root it is \(E\).  As a premise \(sc=sc\)
of form (2), \(sc\) occurs in \(S\).  It is not a side premise of form (1),
by Step~5.  As a top equation of a spine, it is the top of a fragment of
Step~2, kept in Step~4 only if \(a\) occurs in \(E\) or \(S\), or the new
axiom \(a_1=a_1\) of Step~4; deletions in Step~5 keep the top equation of a
spine, and a spine all of whose inferences are deleted leaves its top
equation in place of its bottom, which by Step~3 is the root or a premise of
form (2).  So every axiom satisfies the subterm property, and \(T_S\) has the
subterm property.
\end{proof}
\end{annexitem}

\begin{annexitem}{annex:p02-w02-p02-028-st-prop-3-5}{thm:p02-w02-p02-028-st-prop-3-5}{Obtaining the weak subterm property}
\begin{proof}
Let \(n=\operatorname{lh}(T)\).  Call an inference of form (2) with premise
\(sa=sb\) \emph{good} if \(sa\) and \(sb\) occur in \(S\), and \emph{bad}
otherwise; call an equation an \emph{\(S\)-equation} if both its sides occur
in \(S\).  The conclusion of a good inference of form (2) is an
\(S\)-equation, since \(a\) and \(b\) are subterms of \(sa\) and \(sb\).
For a bad inference \(R\) of \(T\) let \(U_0(R)\) be the subcomputation of
\(T\) ending in its premise and \(p(R)=\operatorname{lh}(U_0(R))\).

\emph{Plan.}  We treat the bad inferences of \(T\) one at a time, as the
chapter does, each time choosing one with no untreated bad inference above
it.  Treating \(R\), with premise \(sa=sb\) and conclusion \(a=b\), replaces
the current subcomputation \(U\) of \(sa=sb\) together with \(R\) by a
singular computation \(N_R\) of \(a=b\) from assumptions of \(U\) whose new
inferences of form (2) are good; the rest of the computation is unchanged.
So throughout, (V) every inference of form (2) of the current computation is
good or is a bad inference of \(T\) not yet treated, with its premise and
conclusion unchanged.  When all are treated, the result \(T_W\) is a
singular computation of \(E\) from \(S\) with the weak subterm property.

\emph{Stage for \(R\).}  Let \(F\) be the \(H\) fragment of \(U\) whose
bottom is \(sa=sb\), or the single equation \(sa=sb\) if that is not the
conclusion of an inference of form (1), and let \(i(F)\) be its number of
inferences.  By the choice of \(R\) and (V), every inference of form (2) in
\(U\) is good, so every leaf of \(F\) is an axiom, an assumption, or the
conclusion of a good inference: each leaf that is not an axiom is an
\(S\)-equation.  This is the induction hypothesis of the chapter: all terms
of the side premises used below occur in \(S\).  If \(F\) is a single
equation, it is an axiom \(sa=sa\), since otherwise \(R\) would be good; then
\(N_R\) is the axiom \(a=a\).  Otherwise:

(a) \emph{Normalizing \(F\).}  Apply, inside \(U\), the inverses of shifts to
pairs of inferences of \(F\) nested to the right, as in part~2 of the proof
of Theorem~\ref{thm:p02-w02-p02-028-st-normal-forms} (each increases the
leftness of \(U\), so this stops); then Step~2 of the proof of
Proposition~\ref{thm:p02-w02-p02-028-st-prop-3-4}; then delete, as in its
Step~5, the inferences of \(F\) whose side premise is an axiom, say \(d\) of
them.  Moves keep the leaves of \(F\), the subcomputations above them, the
bottom \(sa=sb\) and the number \(i(F)\).  Now \(F\) is a spine with top
equation \(a_p=b_q\), left replacements and then right replacements as in
that Step~2, \(i(F)-d\) inferences, side premises that are \(S\)-equations,
and bottom \(a_1=b_1\), which is \(sa=sb\).

(b) \emph{A chain on the right.}  Rebuild \(F\) as a spine that starts from
the axiom \(sa=sa\) and has only right replacements.  If the top is an axiom
\(a_p=a_p\), do as in Step~4 of the proof of
Proposition~\ref{thm:p02-w02-p02-028-st-prop-3-4}.  If not, the top is an
\(S\)-equation; reverse the left replacements from \(a_1=a_1\) to
\(a_1=a_p\) as in that Step~4, then pass to \(a_1=b_q\) by replacing \(a_p\)
by \(b_q\) at position \(2\) with the old top as side premise, then continue
with the right replacements.  This is the chapter's insertion of the
inference from \(a_1=a_p\) and \(a_p=b_q\) to \(a_1=b_q\).  In the chapter's
notation the result is a spine
\[
 sa=t_0,\quad sa=t_1,\quad\ldots,\quad sa=t_{m+1},
\]
\(t_0\) being \(sa\) and \(t_{m+1}\) being \(sb\), whose step from \(sa=t_k\)
to \(sa=t_{k+1}\) is a replacement inside \(t_k\) with an \(S\)-equation
\(E_k\) as side premise.  Every non-axiom leaf of \(F\) is used as exactly
one \(E_k\), with its subcomputation, and \(m+1\leq i(F)-d+1\leq i(F)+1\).

(c) \emph{Case (ii): blocks without \(s\).}  Call a term an \(s\)-term if it
has the form \(sr\).  \(t_0\) and \(t_{m+1}\) are \(s\)-terms.  Take a maximal
block \(t_i,\ldots,t_{j-1}\) (\(1\leq i<j\leq m+1\)) of terms that are not
\(s\)-terms; then \(t_{i-1}\) and \(t_j\) are \(s\)-terms.  The step from
\(t_{i-1}\) to \(t_i\) changes the outer symbol, so it replaces the whole of
\(t_{i-1}\): \(E_{i-1}\) is \(t_{i-1}\doteq t_i\).  Likewise \(E_{j-1}\) is
\(t_{j-1}\doteq t_j\).  So \(t_{i-1}\) and \(t_j\) occur in \(S\).  Replace
the steps from \(sa=t_{i-1}\) to \(sa=t_j\) by the spine that starts from
\(E_{i-1}\) and replaces inside its side \(t_i\) successively with
\(E_i,\ldots,E_{j-1}\), at the positions those steps used, reaching
\(t_{i-1}\doteq t_j\) (in the orientation of \(E_{i-1}\)), followed by the
single step from \(sa=t_{i-1}\) to \(sa=t_j\) that replaces \(t_{i-1}\) by
\(t_j\) at position \(2\) with side premise \(t_{i-1}\doteq t_j\).  Every new
inference is a replacement at one position, hence singular, and
\(t_{i-1}\doteq t_j\) is an \(S\)-equation.  After all blocks are treated, the
main spine has only \(s\)-terms on the right and each of its steps has an
\(S\)-equation as side premise.

(d) \emph{Case (i): removing \(s\).}  Write the right sides of the main spine
as \(sr_0,\ldots,sr_{m'+1}\), with \(r_0\) being \(a\) and \(r_{m'+1}\) being
\(b\), and replace each \(sa=sr_k\) by \(a=r_k\).  A step that replaces inside
the argument of \(s\), at a position \(21\sigma\), becomes the replacement
at \(2\sigma\) in \(a=r_k\) with the same side premise.  A step at position
\(2\) has side premise \(sr_k\doteq sr_{k+1}\), an \(S\)-equation; it becomes
the inference of form (2) from that side premise to \(r_k\doteq r_{k+1}\),
which is good, followed by the replacement of \(r_k\) by \(r_{k+1}\) at
position \(2\).  The result \(N_R\) is a singular computation of \(a=b\) that
starts from the axiom \(a=a\).  Its assumptions are among those of \(U\); its
new inferences of form (2) are good, and its other inferences of form (2) lie
in reused subcomputations of leaves of \(F\), where they are good.  So
replacing \(U\) and \(R\) by \(N_R\) keeps (V), with \(R\) treated, and leaves
the rest of the computation a correct computation, since \(N_R\) ends in
\(a=b\).

\emph{Counting one stage.}  Let \(H\) be the sum, over the non-axiom leaves
of \(F\), of the length of the subcomputation of the leaf minus one.  \(F\)
has \(i(F)\) conclusions and \(i(F)+1\) leaves, so
\(\operatorname{lh}(U)+1=H+2i(F)+2\).  Apart from the reused leaves and their
subcomputations, which contribute \(H+m+1\), \(N_R\) consists of the axiom
\(a=a\) and, for each step of (b), its conclusion and at most one conclusion
of form (2); a block of (c) of \(j-i+1\geq2\) steps instead yields \(j-i\)
conclusions on its own spine, one conclusion of form (2) and one main
conclusion, that is \(j-i+2\leq2(j-i+1)\).  So \(N_R\) has at most
\(2(m+1)+1\) new equations and
\(\operatorname{lh}(N_R)\leq H+3(m+1)+1\).  As \(m\leq i(F)\),
\[
 \operatorname{lh}(N_R)-(\operatorname{lh}(U)+1)\leq3m+4-2i(F)-2\leq i(F)+2 .
\]
When \(F\) is a single axiom the length decreases by one.

\emph{Bounding \(i(F)\).}  We show by induction over the stages that
\(i(F)\) at the stage of \(R\) is at most the number of inferences of
\(U_0(R)\).  Earlier stages replaced subcomputations ending in other bad
inferences; such an inference is of form (2), so it lies above a leaf of the
original fragment of \(sa=sb\) in \(T\) or outside \(U_0(R)\), and the
inferences of that original fragment are unchanged.  If a treated \(R'\) had
its conclusion as a leaf of that fragment, the main spine of \(N_{R'}\) now
joins \(F\) above that leaf; its side premises are leaves (non-axiom leaves
of the old fragment, or conclusions of form (2)), so nothing else joins.  That
spine has \(m'+1\leq i(F_{R'})+1\) inferences, at most the number of
inferences of \(U_0(R')\) plus one, by the induction hypothesis.  The
original fragment's own inferences, and for each such \(R'\) the inference
\(R'\) and the inferences of \(U_0(R')\), are distinct inferences of
\(U_0(R)\).  So \(i(F)\) is at most the number of inferences of \(U_0(R)\),
which is at most \(p(R)-1\), since each inference has its own conclusion and
\(U_0(R)\) has a leaf.

\emph{The total.}  With \(p=p(R)\): \(F\) has \(2i(F)+1\leq2p-1\) equation
occurrences, the main spine of \(N_R\) passed on to the next fragment has
\(m+2\leq p+1\) equations, \(N_R\) has at most \(2m+3\leq2p+1\) new equations,
and the conclusion of \(R\) is deleted; these are within the chapter's figures
\(2p\), \(p+2\) and \(2p+2\).  Each stage adds at most \(i(F)+2\leq p(R)+1\)
to the length, and \(p(R)\leq n-1\) because \(U_0(R)\) omits the conclusion
of \(R\).  The number of stages is at most the number of inferences of form
(2) of \(T\), which is at most \(n-1\) when it is positive.  Hence
\[
 \operatorname{lh}(T_W)\leq n+(n-1)n=n^2\leq n^2+n+1,
\]
and \(T_W=T\) when \(T\) has no bad inference.
\end{proof}
\end{annexitem}

\begin{annexitem}{annex:p02-w02-p02-028-st-bound-b}{thm:p02-w02-p02-028-st-bound-b}{The bound for question (b)}
\begin{proof}
Condition (i) of [1.7] says that \(T\) uses no rule (3) or (\(4_n\)).  Under
condition (ii), an inference of form (1) whose rule variable does not occur
in its equation has conclusion equal to its main premise; replacing the
subcomputation ending in it by that of its main premise shortens \(T\) and
leaves a computation of the same equation from the same assumptions.  After
all such deletions every inference of form (1) has its variable exactly
once, so the computation is singular, and its length is at most \(p\).
Proposition~\ref{thm:p02-w02-p02-028-st-prop-3-5} gives a singular \(T_W\)
with the weak subterm property, of length at most \(p^2\) by the count in
its proof; Proposition~\ref{thm:p02-w02-p02-028-st-prop-3-4} gives a
singular \(T_S\) with the subterm property and no greater length; and
Proposition~\ref{thm:p02-w02-p02-028-st-prop-3-3-2} with
\(M=\operatorname{lh}\) gives an \(\operatorname{lh}\)-direct \(T_M\) to which
\(T_S\) reduces, of the same length by
Lemma~\ref{thm:p02-w02-p02-028-st-moves-basic}.  All three are computations
of \(E\) from \(S\).
\end{proof}
\end{annexitem}

\begin{annexitem}{annex:p02-w02-p02-028-st-semigroup}{thm:p02-w02-p02-028-st-semigroup}{Cancellation-semigroup rules}
\begin{proof}
(2) implies (1): in a cancellation semigroup satisfying the \(R_i\), every
axiom is true, substitution of equals preserves truth by the term-value
argument (T1) (with the product in place of the function variables), and the
cancellation rules preserve truth by cancellativity.

(1) implies (2): let \(D\) be the set of derivable equations and write
\(a\sim b\) for \(a=b\in D\).  Reflexivity: the identity axioms.  Symmetry:
from \(a=b\), rule (1) at \(w\) in \(w=a\), with premise \(a=a\), gives
\(b=a\).  Transitivity: from \(a=b\) and \(b=c\), rule (1) at \(w\) in
\(a=w\) gives \(a=c\).  Compatibility: from \(a=b\), rule (1) at \(w\) in
\((ca)=(cw)\), with premise the axiom \((ca)=(ca)\), gives \((ca)=(cb)\), and
similarly \((ac)=(bc)\).  So \(\sim\) is a congruence and the set \(Q\) of
classes carries the product \([a][b]=[(ab)]\).  It is associative by the
associativity axioms and cancellative by the cancellation rules: if
\([(ab)]=[(ac)]\) then \((ab)=(ac)\in D\), so \(b=c\in D\); similarly on the
right.  Interpreting each generator \(g\) by \([g]\), every word \(t\)
denotes \([t]\), by induction, so each \(R_i\), being in \(D\), is true.  If
\(w_1=w_2\notin D\), then \([w_1]\neq[w_2]\), so \(w_1=w_2\) fails in the
cancellation semigroup \(Q\), contrary to (1).
\end{proof}
\end{annexitem}

\begin{annexitem}{annex:p02-w02-p02-028-st-horn}{thm:p02-w02-p02-028-st-horn}{Horn derivations for the class \(L\)}
\begin{proof}
Since \(\exists\bar xG\) is a sentence, the \(B_k\) are ground atoms and the
variables of the \(A_i\) and \(I_j\) are among \(\bar x\).  If the signature
has no constant, add one.  By
Theorem~\ref{thm:p02-w02-p02-028-existential-validity} there are ground
substitutions \(\theta_1,\ldots,\theta_r\) with
\(G\theta_1\lor\cdots\lor G\theta_r\) propositionally valid.  Write
\(G\theta_l=(X_l\land Y)\to B_{m+1}\), with
\(X_l=\bigwedge_i\theta_lA_i\land\bigwedge_j\theta_lI_j\) and
\(Y=\bigwedge_kB_k\).  An atom valuation falsifies
\(\bigvee_l\bigl((X_l\land Y)\to B_{m+1}\bigr)\) iff it makes every \(X_l\)
and \(Y\) true and \(B_{m+1}\) false.  So every valuation making all
\(\theta_lA_i\), \(\theta_lI_j\) and \(B_k\) true makes \(B_{m+1}\) true.

Let \(D\) be the least set of ground atoms containing every \(\theta_lA_i\)
and every \(B_k\) and closed under the rules
\(\theta_lP^j_1,\ldots,\theta_lP^j_{r(j)}\,/\,\theta_lC^j\); it is the union of
the stages \(D_0\subseteq D_1\subseteq\cdots\), where \(D_{q+1}\) adds the
conclusions of rules whose premises are in \(D_q\), so every member of \(D\)
has a finite derivation by these rules.  The characteristic valuation of
\(D\) makes each \(\theta_lA_i\) and \(B_k\) true and each \(\theta_lI_j\)
true: if all antecedents \(\theta_lP^j_i\) are in \(D\), so is
\(\theta_lC^j\).  Hence it makes \(B_{m+1}\) true, that is,
\(B_{m+1}\in D\), and \(B_{m+1}\) is derivable.
\end{proof}
\end{annexitem}

\begin{annexitem}{annex:p02-w02-p02-028-st-gform}{thm:p02-w02-p02-028-st-gform}{Instances of \(G\)-forms}
\begin{proof}
For a map \(\eta\) from variables to terms and a quantifier-free formula
\(\chi\), write \(\eta\chi\) for the result of replacing every occurrence of
each variable \(u\) in \(\chi\) by \(\eta(u)\), simultaneously.  Three facts
hold by induction on terms and then on the connectives of \(\chi\).
\begin{enumerate}
  \item[(S1)] \(\chi[t/v]=\eta\chi\) for the map \(\eta\) sending \(v\) to
        \(t\) and fixing every other variable: on a quantifier-free formula
        the capture-avoiding substitution substitutes in every atomic term,
        and the term clause replaces exactly the occurrences of \(v\).
  \item[(S2)] \(\eta_2(\eta_1\chi)=(\eta_2\circ\eta_1)\chi\), where
        \((\eta_2\circ\eta_1)(u)\) is the result of replacing each variable
        \(w\) in the term \(\eta_1(u)\) by \(\eta_2(w)\); and \(\eta\chi\)
        depends only on the values of \(\eta\) at the variables of \(\chi\).
  \item[(S3)] If \(\eta\) fixes every member of a set \(P\) of variables and
        \(P\subseteq\operatorname{Var}(\chi)\), then
        \(P\subseteq\operatorname{Var}(\eta\chi)\), since an occurrence of
        \(u\in P\) is replaced by \(u\).
\end{enumerate}

\emph{Claim.}  Let \(\chi\) be quantifier-free, \(P\) a set of variables with
\(P\subseteq\operatorname{Var}(\chi)\), \(w_1,\ldots,w_r\) (\(r\geq0\))
variables none of which lies in \(P\), \(v\notin P\) a variable and \(t\) a
term.  Then
\[
  (\exists w_1\cdots\exists w_r\,\chi)[t/v]
  =\exists w'_1\cdots\exists w'_r\,\eta\chi
\]
for some variables \(w'_1,\ldots,w'_r\), none in \(P\), and some map
\(\eta\) that fixes every member of \(P\).

The claim is proved by induction on \(r\), for all \(\chi,P,w_i,v,t\) at once.
If \(r=0\), (S1) gives the claim with \(\eta\) sending \(v\) to \(t\); it
fixes \(P\) because \(v\notin P\).  Let \(r>0\), and write
\(\psi_1=\exists w_2\cdots\exists w_r\,\chi\).  The definition of
capture-avoiding substitution
(Definition~\ref{def:p02-fol-capture-avoiding-substitution}) has three
clauses for \((\exists w_1\psi_1)[t/v]\).
\begin{enumerate}
  \item If \(w_1=v\), the result is \(\exists w_1\psi_1\) itself: take
        \(w'_i=w_i\) and \(\eta\) the identity.
  \item If \(w_1\neq v\) and \(w_1\notin\operatorname{Var}(t)\), the result
        is \(\exists w_1(\psi_1[t/v])\).  The induction hypothesis for
        \(\psi_1\), which has \(r-1\) quantifiers, gives
        \(\psi_1[t/v]=\exists w'_2\cdots\exists w'_r\,\eta\chi\) with the
        required properties; put \(w'_1=w_1\notin P\).
  \item If \(w_1\neq v\) and \(w_1\in\operatorname{Var}(t)\), the result is
        \(\exists z((\psi_1[z/w_1])[t/v])\), where \(z\) is the least
        variable outside
        \(\operatorname{Var}(\psi_1)\cup\operatorname{Var}(t)\cup\{v,w_1\}\).
        The induction hypothesis for \(\psi_1\), the variable \(w_1\notin P\)
        and the term \(z\) gives
        \(\psi_1[z/w_1]=\exists w''_2\cdots\exists w''_r\,\eta_1\chi\) with no
        \(w''_i\) in \(P\) and \(\eta_1\) fixing \(P\).  By (S3),
        \(P\subseteq\operatorname{Var}(\eta_1\chi)\), so the induction
        hypothesis applies again, to this formula, the variable \(v\) and the
        term \(t\), and gives
        \((\psi_1[z/w_1])[t/v]=\exists w'_2\cdots\exists w'_r\,\eta_2(\eta_1\chi)\)
        with no \(w'_i\) in \(P\) and \(\eta_2\) fixing \(P\).  By (S2) the
        matrix is \((\eta_2\circ\eta_1)\chi\), and \(\eta_2\circ\eta_1\)
        fixes \(P\).  Put \(w'_1=z\).  Every member of \(P\) occurs in
        \(\chi\) and is none of \(w_2,\ldots,w_r\), so it occurs free in
        \(\psi_1\) and lies in \(\operatorname{Var}(\psi_1)\); hence
        \(z\notin P\).
\end{enumerate}

\emph{The lemma.}  Let \(\exists w_1\psi\) be a \(G\)-form with
\(\psi=\exists w_2\cdots\exists w_s\,\theta G\), and let
\(P=\operatorname{par}(G)\).  The substitution \(\theta\) moves only
\(x_1,\ldots,x_n\), none of which is in \(P\), so it fixes \(P\); since
\(P\subseteq\operatorname{Var}(G)\), (S3) gives
\(P\subseteq\operatorname{Var}(\theta G)\).  Also \(w_1\notin P\), and no
\(w_i\) is in \(P\).  The claim, applied to \(\chi=\theta G\), \(v=w_1\) and
\(t\), gives
\(\psi[t/w_1]=\exists w'_2\cdots\exists w'_s\,\eta(\theta G)\) with no
\(w'_i\) in \(P\) and \(\eta\) fixing \(P\).  By (S2),
\(\eta(\theta G)=(\eta\circ\theta)G\).  On the variables of \(G\) the map
\(\eta\circ\theta\) agrees with the substitution \(\theta'\) given by
\(\theta'(x_i)=\eta(\theta(x_i))\): at \(x_i\) by definition, and at
\(u\in P\) because \((\eta\circ\theta)(u)=\eta(u)=u=\theta'(u)\).  By the
second part of (S2), \(\eta(\theta G)=\theta'G\), so \(\psi[t/w_1]\) is a
\(G\)-form with matrix substitution \(\theta'\).
\end{proof}
\end{annexitem}

\begin{annexitem}{annex:p02-w02-p02-028-st-lk-translation}{thm:p02-w02-p02-028-st-lk-translation}{From canonical LK to the defined-negation calculus}
\begin{proof}
Two facts about \(d\) come first.

(D1) \(\operatorname{Var}(d\varphi)=\operatorname{Var}(\varphi)\) and
\(\FV(d\varphi)=\FV(\varphi)\) for every formula \(\varphi\), by induction on
\(c(\varphi)\): the only clause of \(d\) that changes the shape of a formula
replaces \(\neg\psi\) by \(d\psi\to\bot\), and \(\bot\) contains no variable.

(D2) \(d(\varphi[t/x])=(d\varphi)[t/x]\) for every formula \(\varphi\), term
\(t\) and variable \(x\).  The proof is by strong induction on \(c(\varphi)\).
For an atomic formula or \(\bot\), \(d\) is the identity on both sides, since
substitution turns an atomic formula into an atomic formula and leaves
\(\bot\) unchanged.  For \(\neg\psi\):
\(d((\neg\psi)[t/x])=d(\neg(\psi[t/x]))=d(\psi[t/x])\to\bot
=(d\psi)[t/x]\to\bot=(d\psi\to\bot)[t/x]\), by the induction hypothesis for
\(\psi\) and because substitution commutes with \(\to\) and fixes \(\bot\).
The binary connectives are treated in the same way.  For \(Qy\,\psi\) we have
\(d(Qy\,\psi)=Qy\,d\psi\).  Which clause of the substitution definition
applies depends only on \(y\), \(x\) and \(\operatorname{Var}(t)\), so the
same clause applies on both sides; in the renaming clause the new variable is
the least one outside
\(\operatorname{Var}(\psi)\cup\operatorname{Var}(t)\cup\{x,y\}\), and by
(D1) this set equals
\(\operatorname{Var}(d\psi)\cup\operatorname{Var}(t)\cup\{x,y\}\), so the
same variable \(z\) is chosen.  In the first clause both sides are
\(Qy\,d\psi\).  In the second,
\(d(Qy(\psi[t/x]))=Qy\,d(\psi[t/x])=Qy((d\psi)[t/x])\).  In the third,
\[
 d\bigl(Qz((\psi[z/y])[t/x])\bigr)
 =Qz\bigl((d(\psi[z/y]))[t/x]\bigr)
 =Qz\bigl(((d\psi)[z/y])[t/x]\bigr)
 =(Qy\,d\psi)[t/x],
\]
where the first step uses the induction hypothesis at \(\psi[z/y]\), whose
complexity is \(c(\psi)<c(Qy\,\psi)\) by
Lemma~\ref{lem:substitution-complexity}, and the second uses it at \(\psi\).

Now induct on a derivation on \texttt{sheet:lk}.  For each inference with
conclusion \(\Gamma\Rightarrow\Delta\) we derive
\(d\Gamma\Rightarrow d\Delta\) on \texttt{sheet:lk-defneg} from the
translated premises.
\begin{enumerate}
  \item \(\mathrm{id}\), \(\Gamma,\varphi\Rightarrow\varphi,\Delta\): the
        initial sequent \(d\varphi\Rightarrow d\varphi\), followed by left
        weakening by each member of \(d\Gamma\) and right weakening by each
        member of \(d\Delta\).
  \item \(\bot L\), \(\Gamma,\bot\Rightarrow\Delta\): since \(d\bot=\bot\),
        this is the initial sequent \(\bot,d\Gamma\Rightarrow d\Delta\).
  \item Weakening and contraction on either side: the same rule.
  \item \(\neg L\): from \(d\Gamma\Rightarrow d\Delta,d\varphi\) and the
        initial sequent \(\bot\Rightarrow\) (that is, \(\bot L\) with empty
        contexts), \(\to L\) gives
        \(d\varphi\to\bot,d\Gamma\Rightarrow d\Delta\), which is
        \(d(\neg\varphi),d\Gamma\Rightarrow d\Delta\).
  \item \(\neg R\): from \(d\varphi,d\Gamma\Rightarrow d\Delta\), right
        weakening by \(\bot\) and then \(\to R\) give
        \(d\Gamma\Rightarrow d\Delta,d\varphi\to\bot\).
  \item \(\land L\), \(\land R\), \(\lor L\), \(\lor R\), \(\to R\): the rule
        of the same name, since \(d\) commutes with these connectives and the
        context conventions agree.
  \item \(\to L\): the canonical rule shares its contexts.  From
        \(d\Gamma\Rightarrow d\Delta,d\varphi\) and
        \(d\psi,d\Gamma\Rightarrow d\Delta\), the split rule gives
        \(d\varphi\to d\psi,d\Gamma,d\Gamma\Rightarrow d\Delta,d\Delta\), and
        left and right contraction on each repeated member give
        \(d(\varphi\to\psi),d\Gamma\Rightarrow d\Delta\).
  \item Cut and Mix: both sheets split the contexts in the same way, and
        Mix on \texttt{sheet:lk-defneg} allows further occurrences of the mix
        formula in the side contexts, so the translated inference is an
        instance of the rule of the same name.
  \item \(\forall L\) and \(\exists R\) with term \(t\): by (D2) the
        translated premise contains \((d\varphi)[t/x]\), and the rule of the
        same name applies.
  \item \(\forall R\) and \(\exists L\) with eigenvariable \(y\): by (D2) the
        translated premise contains \((d\varphi)[y/x]\).  On
        \texttt{sheet:lk}, \(y\) is free in no formula of the conclusion; by
        (D1) it is free in no formula of the translated conclusion, which is
        the freshness condition of \texttt{sheet:lk-defneg}.
\end{enumerate}
These cases cover every rule of \texttt{sheet:lk}.
\end{proof}
\end{annexitem}

\begin{annexitem}{annex:p02-w02-p02-028-st-extraction}{thm:p02-w02-p02-028-st-extraction}{Horn derivations from cut-free derivations}
\begin{proof}
Four preliminary facts.

(E1) \emph{Monotonicity.}  If \(\Theta\subseteq\Theta'\) and every member of
\(\Gamma\) is a member of \(\Gamma'\), then every axiom and every rule of
\(\mathsf H(\Theta;\Gamma)\) belongs to \(\mathsf H(\Theta';\Gamma')\), so
every derivation in the first is a derivation in the second.  In particular
\(\mathsf H(\Theta;\Gamma)\) depends on \(\Gamma\) only through the set of
its members.

(E2) \emph{Shapes.}  A left formula is a prime formula, a conjunction
\(\theta E'\), or an implication whose antecedent is prime (a tail).  A right
formula is a prime formula, a formula beginning with \(\exists\) (a
\(G\)-form with \(s\geq1\)), or an implication \(\theta G=\theta E\to
B_{m+1}\) (a \(G\)-form with \(s=0\)).  A formula that is both a left and a
right formula is prime: a \(G\)-form is not a conjunction, and if
\(\theta G\) were a tail of some \(\theta'I_j\), then \(q\geq1\) and
\(\theta E\) would be prime, whereas \(\theta E\) is a conjunction when
\(p+q+m\geq2\) and is the implication \(\theta I_1\) when \(p+q+m=1\) and
\(q\geq1\).  No left or right formula is \(\bot\) or a disjunction or begins
with \(\forall\), and no left formula begins with \(\exists\).

(E3) \emph{Parts of the antecedent.}  For a substitution \(\theta\) and a
subformula \(E'\) of \(E\), the formula \(\theta E'\) is a left formula: a
conjunction \(\theta E'\) if \(E'\) is a conjunction, the prime formula
\(\theta A_i\) or \(B_k\) if \(E'\) is the leaf \(A_i\) or \(B_k\), and the
first tail \(\theta I_j\) if \(E'\) is the leaf \(I_j\).  In every case the
axioms and rules that \(\theta E'\) contributes as a member of \(\Gamma\) are
exactly the axioms \(\theta A_i\) and \(B_k\) and the rules
\(\theta P^j_1\cdots\theta P^j_{r(j)}/\theta C^j\) for the leaves
\(A_i\), \(B_k\), \(I_j\) of \(E'\).

(E4) \emph{Targets.}  If every prime member of \(\Delta'\) is a member of
\(\Delta\), and \(\Delta\) has a \(G\)-form member whenever \(\Delta'\) does,
then every target of \(\Delta'\) is a target of \(\Delta\).

The proof is by induction on the height of \(\mathcal D\), by cases on its
last inference.  Cut and Mix do not occur.  In each case the premises again
have left formulas in the antecedent and right formulas in the succedent; this
is checked case by case below.
\begin{enumerate}
  \item \emph{Id}, \(\varphi\Rightarrow\varphi\).  By (E2) \(\varphi\) is
        prime.  It is an axiom of \(\mathsf H(\varnothing;\{\varphi\})\) and
        a target of \(\{\varphi\}\); take \(\Theta=\varnothing\).
  \item \emph{\(\bot L\)}.  Its antecedent contains \(\bot\), which by (E2)
        is not a left formula; the case does not arise.
  \item \emph{Left weakening and left contraction.}  The premise antecedent
        has only members of \(\Gamma\) as members and the same succedent.  The
        induction hypothesis gives \(\Theta\) and a target of \(\Delta\)
        derivable in the premise system, hence in
        \(\mathsf H(\Theta;\Gamma)\) by (E1).
  \item \emph{Right weakening and right contraction.}  The premise succedent
        has only members of \(\Delta\) as members, so by (E4) its targets are
        targets of \(\Delta\); the antecedent is unchanged.
  \item \emph{\(\land L\)} with principal formula \(\varphi\land\psi\in\Gamma\).
        By (E2) it is \(\theta E'\) with \(E'=E_1\land E_2\), so
        \(\varphi=\theta E_1\), \(\psi=\theta E_2\), and these are left
        formulas by (E3).  The induction hypothesis gives \(\Theta\) and a
        target of \(\Delta\) derivable in
        \(\mathsf H(\Theta;\theta E_1,\theta E_2,\Gamma')\), where
        \(\Gamma=\theta E',\Gamma'\).  By (E3) the contributions of
        \(\theta E_1\) and \(\theta E_2\) are those of the leaves of \(E_1\)
        and \(E_2\), which are leaves of \(E'\), and so are contributed by
        \(\theta E'\in\Gamma\).  So every axiom and rule of the premise system
        belongs to \(\mathsf H(\Theta;\Gamma)\).
  \item \emph{\(\land R\), \(\lor L\), \(\lor R\), \(\forall L\),
        \(\forall R\), \(\exists L\).}  By (E2) there is no right
        conjunction, no disjunction, no formula beginning with \(\forall\),
        and no left formula beginning with \(\exists\); these cases do not
        arise.
  \item \emph{\(\to R\)} with principal formula in \(\Delta\).  By (E2) it is
        some \(\theta G=\theta E\to B_{m+1}\), recalling
        \(\theta B_{m+1}=B_{m+1}\).  The premise is
        \(\theta E,\Gamma\Rightarrow\Delta',B_{m+1}\) with
        \(\Delta=\Delta',\theta G\); \(\theta E\) is a left formula by (E3)
        and \(B_{m+1}\) is prime.  The induction hypothesis gives
        \(\Theta'\) and a target \(\beta\) of \(\Delta',B_{m+1}\) derivable in
        \(\mathsf H(\Theta';\theta E,\Gamma)\).  Put
        \(\Theta=\Theta'\cup\{\theta\}\).  By (E3), \(\theta E\) contributes
        the axioms \(\theta A_i\), \(B_k\) and the rules for \(\theta I_j\),
        all of which belong to \(\mathsf H(\Theta;\Gamma)\) because
        \(\theta\in\Theta\); with (E1), \(\beta\) is derivable in
        \(\mathsf H(\Theta;\Gamma)\).  If \(\beta\) is a prime member of
        \(\Delta'\), it is a target of \(\Delta\).  Otherwise
        \(\beta=B_{m+1}\), and \(B_{m+1}\) is a target of \(\Delta\) because
        \(\theta G\in\Delta\) is a \(G\)-form.  The new member \(\theta\) of
        \(\Theta\) is the matrix substitution of the principal formula
        \(\theta G\).
  \item \emph{\(\to L\)} with principal formula in \(\Gamma\).  By (E2) it is
        an \(l\)-th tail \(\theta(P^j_l\to R_l)\) of some \(\theta I_j\),
        where \(R_l=P^j_{l+1}\to(\cdots(P^j_{r(j)}\to C^j)\cdots)\) if
        \(l<r(j)\) and \(R_l=C^j\) if \(l=r(j)\).  The premises are
        \(\Gamma_1\Rightarrow\Delta_1,\theta P^j_l\) and
        \(\theta R_l,\Gamma_2\Rightarrow\Delta_2\), with
        \(\Gamma=\theta(P^j_l\to R_l),\Gamma_1,\Gamma_2\) and
        \(\Delta=\Delta_1,\Delta_2\).  Here \(\theta P^j_l\) is prime, and
        \(\theta R_l\) is the \((l+1)\)-th tail or the prime formula
        \(\theta C^j\); so both premises have the required form.  The
        induction hypothesis gives \(\Theta_1\) and a target \(\beta_1\) of
        \(\Delta_1,\theta P^j_l\) derivable in
        \(\mathsf H(\Theta_1;\Gamma_1)\), and \(\Theta_2\) and a target
        \(\beta_2\) of \(\Delta_2\) derivable in
        \(\mathsf H(\Theta_2;\theta R_l,\Gamma_2)\).  Put
        \(\Theta=\Theta_1\cup\Theta_2\) and \(\mathsf H=\mathsf H(\Theta;\Gamma)\).

        If \(\beta_1\) is a target of \(\Delta_1\), it is a target of
        \(\Delta\) by (E4) and is derivable in \(\mathsf H\) by (E1).
        Otherwise \(\beta_1=\theta P^j_l\), which is derivable in \(\mathsf H\)
        by (E1); let \(\delta_0\) be such a derivation.  Since \(\beta_2\) is a
        target of \(\Delta\) by (E4), it suffices to derive it in
        \(\mathsf H\).  Every axiom and rule of
        \(\mathsf H(\Theta_2;\theta R_l,\Gamma_2)\) belongs to \(\mathsf H\)
        except possibly the one contributed by \(\theta R_l\): the axiom
        \(\theta C^j\) if \(l=r(j)\), the rule
        \(\theta P^j_{l+1}\cdots\theta P^j_{r(j)}/\theta C^j\) if \(l<r(j)\).
        The tail \(\theta(P^j_l\to R_l)\in\Gamma\) contributes to \(\mathsf H\)
        the rule \(\theta P^j_l\cdots\theta P^j_{r(j)}/\theta C^j\).  Let
        \(\delta\) derive \(\beta_2\) in
        \(\mathsf H(\Theta_2;\theta R_l,\Gamma_2)\).  In the sequence
        \(\delta_0\) followed by \(\delta\), every line of \(\delta_0\) is
        justified in \(\mathsf H\); every line of \(\delta\) justified by an
        axiom or rule of \(\mathsf H\) keeps its justification; and a line
        \(\theta C^j\) justified by the contribution of \(\theta R_l\) is the
        conclusion of the rule
        \(\theta P^j_l\cdots\theta P^j_{r(j)}/\theta C^j\) of \(\mathsf H\),
        whose premise \(\theta P^j_l\) occurs in \(\delta_0\) and whose
        remaining premises occur earlier in \(\delta\).  So \(\beta_2\) is
        derivable in \(\mathsf H\).
  \item \emph{\(\exists R\)} with principal formula
        \(\exists w_1\psi\in\Delta\).  By (E2) it is a \(G\)-form with at
        least one quantifier; the premise is
        \(\Gamma\Rightarrow\Delta',\psi[t/w_1]\) with
        \(\Delta=\Delta',\exists w_1\psi\).  By
        Lemma~\ref{thm:p02-w02-p02-028-st-gform}, \(\psi[t/w_1]\) is a
        \(G\)-form, so it is a right formula and is not prime.  The two
        succedents have the same prime members and both contain a \(G\)-form,
        so by (E4) they have the same targets, and the induction hypothesis
        gives the claim.
\end{enumerate}
These cases cover every rule of \texttt{sheet:lk-defneg} other than Cut and
Mix.  In each case the members of \(\Theta\) are inherited from the premises
or, in the case \(\to R\), added as the matrix substitution of the principal
formula, which gives the final statement of the lemma.
\end{proof}
\end{annexitem}

\begin{annexitem}{annex:p02-w02-p02-028-st-horn-cut}{thm:p02-w02-p02-028-st-horn-cut}{The Horn statement of [1.3(ii)] from cut elimination}
\begin{proof}
(1) By Theorem~\ref{thm:p07-classical-cut-elimination}, the given derivation
can be transformed into a derivation of the same literal endsequent
\(\Rightarrow\exists x_1\cdots\exists x_nG\) containing neither Cut nor Mix.
Its antecedent is empty, and its succedent consists of
\(\exists x_1\cdots\exists x_nG\), which is a \(G\)-form with the identity as
matrix substitution: no \(x_i\) lies in \(\operatorname{par}(G)\).
Lemma~\ref{thm:p02-w02-p02-028-st-extraction} gives a finite set
\(\Theta\) and a target of \(\{\exists x_1\cdots\exists x_nG\}\) derivable in
\(\mathsf H(\Theta;\varnothing)\).  That formula is not prime (for \(n=0\) it
is the implication \(G\)), so the only target is \(B_{m+1}\).  The system
\(\mathsf H(\Theta;\varnothing)\) has exactly the assumptions
\(B_1,\ldots,B_m\), the axioms \(\theta A_i\) and the rules
\(\theta P^j_1\cdots\theta P^j_{r(j)}/\theta C^j\) for \(\theta\in\Theta\),
which proves the first claim.  Each of these axioms and rules is among those
of [1.3(ii)] for all substitutions, so the same sequence is a derivation in
that system.

(2) Let \(\Lang_0\) be the finite signature of the symbols occurring in
\(G\).  Every \(\Lang_0\)-structure \(\mathcal N'\) is the reduct of a
structure \(\mathcal N\) for the full signature, by the reduct fact at the
head of this annex; for each assignment \(s\), \(\mathcal N,s\) satisfies
\(\exists x_1\cdots\exists x_nG\) by validity, so \(\mathcal N',s\) does by
the same fact.  Hence the sequent \(\Rightarrow\exists x_1\cdots\exists x_nG\)
is valid under every \(\Lang_0\)-structure and assignment, and
Theorem~\ref{thm:p07-canonical-lk-completeness}, in the case of a finite
signature, for which it uses no form of Choice, gives a derivation of it on
\texttt{sheet:lk} over \(\Lang_0\).  Its formulas and terms are formulas and
terms of the full signature, so it is also a derivation over that signature.
Lemma~\ref{thm:p02-w02-p02-028-st-lk-translation} turns it into a derivation
of \(\Rightarrow d(\exists x_1\cdots\exists x_nG)\) on
\texttt{sheet:lk-defneg}.  By the clauses of
Definition~\ref{def:p07-defined-negation-translation}, \(d\) is the identity
on formulas without \(\neg\), by induction on complexity, and \(G\) contains
no \(\neg\); so the endsequent is \(\Rightarrow\exists x_1\cdots\exists x_nG\)
and item~1 applies.  The reduct fact chooses a single element of a nonempty
domain, so no Choice enters there either.
\end{proof}
\end{annexitem}

\begin{annexitem}{annex:p02-w02-p02-028-st-cycle}{thm:p02-w02-p02-028-st-cycle}{The example \(S_n\)}
\begin{proof}
(1) Under an interpretation satisfying \(S_n\), \(I(x_1)=I(x_n)+1=
I(x_{n-1})+2=\cdots=I(x_1)+n\), which is impossible.

(2) Start from the assumption \(x_1=sx_n\).  From \(x_1=s^qx_{n-q+1}\) and
the assumption \(x_{n-q+1}=sx_{n-q}\), rule (1) at \(w\) in \(x_1=s^qw\)
gives \(x_1=s^{q+1}x_{n-q}\) (indices modulo \(n\)); at \(q=n-1\) this is
\(x_1=s^nx_1\).  Every term has length at most \(n+1=\operatorname{lh}(s^nx_1)\).
Rule (\(4_n\)) then gives \(0=s0\).

(3) Let \(\mathbb Z_n\) be the integers modulo \(n\) with \(0\) and
\(s(v)=v+1\), function variables interpreted arbitrarily.  Successor is a
bijection, so rules (1) and (2) preserve truth there, by the argument of
(T1).  For each \(c\), \(I_c(x_i)=i+c\) satisfies \(S_n\).  Suppose \(T\) is
an \(\operatorname{lh}\)-direct computation of \(0=s0\) from \(S_n\).  The
terms of \(0=s0\) and \(S_n\) have length at most \(2\), so all terms of \(T\)
do.  If \(T\) contained an inference of form (3) or (\(4_m\)), choose one
with no such inference above it; its premise is computed by rules (1) and (2)
from \(S_n\), hence true in \(\mathbb Z_n\) under every \(I_c\).  A premise
\(0=sa\) has \(\operatorname{lh}(a)\leq1\), so \(a\) is \(0\) or a variable;
\(0=s0\) is false as \(n\geq2\), and \(0=sa\) for a variable \(a\) is false
under an \(I_c\) giving \(a\) a value other than \(n-1\).  A premise
\(a=s^ma\) with \(m\geq1\) has \(m+\operatorname{lh}(a)\leq2\), so it is
\(a=sa\), false as \(n\geq2\).  So \(T\) uses only rules (1) and (2), and
\(0=s0\) is true in \(\mathbb Z_n\), which is false for \(n\geq2\).  For
\(n=1\), the inference by rule (\(4_1\)) from \(x_1=sx_1\) to \(0=s0\) has
terms of length at most \(2\).
\end{proof}
\end{annexitem}

\begin{annexitem}{annex:p02-w02-p02-028-st-prop-4-1-1}{thm:p02-w02-p02-028-st-prop-4-1-1}{Substitution instances and directness}
\begin{proof}
\emph{Claim I (instances of computations).}  Let \(T\) be a computation of
\(e\) from \(S'\) by rules (1) and (2), and \(\theta\) a substitution.  Rename
the rule variable of each inference of form (1) to one outside the domain of
\(\theta\) and not occurring in any value of \(\theta\), and apply \(\theta\) to
every equation.  The result \(\theta T\) is a computation of \(\theta e\) from
\(\theta S'\) of the same length.  Axioms \(a=a\) go to axioms and assumptions
to members of \(\theta S'\).  For an inference of form (1) with equation
\(E'\) and variable \(x\), Claim~1 of the proof of
Lemma~\ref{thm:p02-w02-p02-028-st-measure-transfer}, applied to each side,
gives \(\theta(E'(a/x))=(\theta E')(\theta a/x)\) and
\(\theta(E'(b/x))=(\theta E')(\theta b/x)\), so the images form an inference
of form (1) with equation \(\theta E'\).  For form (2),
\(\theta(sa)=s(\theta a)\).  (Singularity need not be preserved.)

\emph{Claim II (singular computations exist).}  Every computation \(T\) of
\(e\) from \(S'\) by rules (1) and (2) can be turned into a singular one of
\(e\) from \(S'\).  By recursion on \(T\): leaves stay; an inference of form
(2) is kept above the transformed premise.  For an inference of form (1) with
equation \(E'\), variable \(x\) occurring \(k\) times in \(E'\), and
transformed premise computations \(T_1'\) of \(E'(a/x)\) and \(T_2'\) of
\(a\doteq b\): if \(k=0\), the conclusion \(E'(b/x)\) is \(E'(a/x)\) and
\(T_1'\) is used.  If \(k\geq1\), let \(E''\) be \(E'\) with its \(k\)
occurrences of \(x\) replaced by distinct new variables \(x_1,\ldots,x_k\),
occurring in none of \(E',a,b\).  Then \(E'(a/x)\) is
\(E''(a/x_1,\ldots,a/x_k)\), and for \(h=1,\ldots,k\) the inference of form
(1) with equation \(E''(b/x_1,\ldots,b/x_{h-1},a/x_{h+1},\ldots,a/x_k)\),
in which \(x_h\) occurs once, and a copy of \(T_2'\) as side premise, leads to
\(E''(b/x_1,\ldots,b/x_h,a/x_{h+1},\ldots,a/x_k)\); the last of these is
\(E'(b/x)\).

\emph{Construction of the instances.}  Let \(f_1,\ldots,f_r\) be the
function variables occurring in \(E\) or \(S\), \(f_l\) of arity \(n_l\).
For each \(l\) and each subset \(J\subseteq\{1,\ldots,n_l\}\) choose a new
function variable \(f_{l,J}\) of arity \(\abs J\), an individual variable
when \(J=\varnothing\), all distinct and occurring nowhere in \(E\) or \(S\).
For each choice \(\kappa=(J_1,\ldots,J_r)\) of subsets let \(\sigma_\kappa\)
send \(f_l\) to \(\lambda x_1\cdots x_{n_l}\,f_{l,J_l}x_{j_1}\cdots x_{j_q}\),
where \(j_1<\cdots<j_q\) list \(J_l\) (to \(\lambda x_1\cdots x_{n_l}\,f_{l,\varnothing}\)
when \(J_l\) is empty), and fix every individual variable.  The pairs
\(E_\kappa=\sigma_\kappa E\), \(S_\kappa=\sigma_\kappa S\), finitely many,
are the instances \(E_1,S_1;\ldots;E_n,S_n\).  Each \(\sigma_\kappa\) is a
guess at which arguments of each \(f_l\) a later substitution will
discard; this is the chapter's anticipation of projections.

\emph{(a)}  If \(E\) is derivable from \(S\), then \(E_\kappa\) is derivable
from \(S_\kappa\), by Claim~I with \(\theta=\sigma_\kappa\).

\emph{(b)}  Let \(\theta\) be any substitution.  For each \(l\) let
\(\theta f_l=\lambda x_1\cdots x_{n_l}\,a_l\), with \(a_l\) the term
\(f_lx_1\cdots x_{n_l}\) when \(f_l\) is not in the domain of \(\theta\);
let \(J_l\) be the set of \(j\) such that \(x_j\) occurs in \(a_l\), and
\(\kappa=(J_1,\ldots,J_r)\).  Let \(\theta_\kappa\) agree with \(\theta\) on
individual variables of \(E\) and \(S\), send \(f_{l,J_l}\) to
\(\lambda x_{j_1}\cdots x_{j_q}\,a_l\) (to \(a_l\) when \(J_l\) is empty,
\(a_l\) then containing no \(x_j\)), and fix every other variable.  Each
bound variable of a value of \(\theta_\kappa\) occurs in its body, so
\(\theta_\kappa\) is non-projecting.  By induction on terms \(t\) of \(E\)
and \(S\), \(\theta_\kappa(\sigma_\kappa t)=\theta t\): for variables, \(0\)
and \(s\) this is immediate; for \(t=f_lt_1\cdots t_{n_l}\),
\(\sigma_\kappa t=f_{l,J_l}(\sigma_\kappa t_{j_1})\cdots(\sigma_\kappa t_{j_q})\),
so
\[
\begin{aligned}
 \theta_\kappa(\sigma_\kappa t)
 &=a_l\bigl(\theta_\kappa\sigma_\kappa t_{j_1}/x_{j_1},\ldots\bigr)
  =a_l\bigl(\theta t_{j_1}/x_{j_1},\ldots\bigr)\\
 &=a_l(\theta t_1/x_1,\ldots,\theta t_{n_l}/x_{n_l})=\theta t,
\end{aligned}
\]
using the induction hypothesis and that \(x_j\) does not occur in \(a_l\) for
\(j\notin J_l\).  Hence \(\theta E=\theta_\kappa E_\kappa\) and
\(\theta S=\theta_\kappa S_\kappa\).

\emph{The proposition.}  Suppose \(E\) is derivable from \(S\).  For each
\(\kappa\), by (a) and Claim~II there is a singular computation of
\(E_\kappa\) from \(S_\kappa\); by
Proposition~\ref{thm:p02-w02-p02-028-st-prop-3-5} one with the weak subterm
property, and by Proposition~\ref{thm:p02-w02-p02-028-st-prop-3-4} a
singular \(T_\kappa\) with the subterm property.  Given \(\theta\), choose
\(\kappa\) and \(\theta_\kappa\) as in (b); (i) and (ii) hold.
\(M^*(c)=M(\theta_\kappa c)\) is a measure by
Lemma~\ref{thm:p02-w02-p02-028-st-measure-transfer}.  By
Proposition~\ref{thm:p02-w02-p02-028-st-prop-3-3-2} (which the printed proof
calls ``Proposition 3.3''), \(T_\kappa\) reduces to an \(M^*\)-direct
\(T_\theta\): (iii).  By Claim~I, \(\theta_\kappa T_\theta\) is a computation
of \(\theta_\kappa E_\kappa=\theta E\) from \(\theta_\kappa S_\kappa=\theta S\).
For (iv), a term occurring in \(\theta_\kappa T_\theta\) is a subterm of a
side \(\theta_\kappa e\), \(e\) a side of an equation of \(T_\theta\), so by
step (B0) of the proof of
Proposition~\ref{thm:p02-w02-p02-028-st-prop-3-3-2} it suffices to bound
\(\theta_\kappa e\).  By \(M^*\)-directness some \(c\) occurring in
\(E_\kappa\) or \(S_\kappa\) has \(M^*(e)\leq M^*(c)\); if \(c\) is a subterm
of a side \(d\) of \(E_\kappa\) or of a member of \(S_\kappa\), then
\(M^*(c)\leq M^*(d)\) by (B0) for the measure \(M^*\).  So
\(M(\theta_\kappa e)\leq M(\theta_\kappa d)\), and \(\theta_\kappa d\) is a
side of \(\theta E\) or of a member of \(\theta S\).  Thus
\(\theta_\kappa T_\theta\) is \(M\)-direct.  Taking \(T_i=T_\kappa\) and
\(\theta_i=\theta_\kappa\) gives the statement.
\end{proof}
\end{annexitem}

\begin{annexitem}{annex:p02-w02-p02-028-st-optimality}{thm:p02-w02-p02-028-st-optimality}{Efficiency of the rules of \(K\)}
\begin{proof}
For each rule \(\rho\in\mathcal R\), \(\bigwedge S_\rho\to E_\rho\) is
analytic, so \(E_\rho\) is derivable from \(S_\rho\) by rules (1) and (2)
(Theorem~\ref{thm:p02-w02-p02-028-st-completeness}(1)).
Proposition~\ref{thm:p02-w02-p02-028-st-prop-4-1-1} with
\(M=\operatorname{lh}\) gives finitely many instances with computations
\(T_{\rho,1},\ldots,T_{\rho,n_\rho}\).  Let \(c\) be the larger of \(1\) and
the greatest length of all these computations, for all \(\rho\in\mathcal R\).

Let \(D\) be as in the statement.  Build the new sequence member by member.
An axiom or a member of \(S\) is copied.  A member \(\theta E_\rho\) obtained
by an application of \(\rho\) from the members of \(\theta S_\rho\), which
occur earlier in \(D\), is replaced as follows.  The proposition gives \(i\),
a non-projecting \(\theta_i\) and \(T_\theta\) such that \(T_{\rho,i}\)
reduces to \(T_\theta\), so that
\(\operatorname{lh}(T_\theta)=\operatorname{lh}(T_{\rho,i})\leq c\)
(Lemma~\ref{thm:p02-w02-p02-028-st-moves-basic}), and such that
\(\theta_iT_\theta\) is an \(\operatorname{lh}\)-direct computation of
\(\theta_iE_{\rho,i}=\theta E_\rho\) from \(\theta_iS_{\rho,i}=\theta S_\rho\),
of the same length by Claim~I of the previous proof.  List the equation
occurrences of \(\theta_iT_\theta\) that are not leaves in members of
\(\theta S_\rho\), each after the premises of the inference concluding it;
the leaves in \(\theta S_\rho\) are not listed again, because they already
stand earlier in the new sequence, as images of earlier members of \(D\).
This adds at most \(c\) equations, the last being \(\theta E_\rho\), and each
is an axiom or follows by rule (1) or (2) from earlier members.  So the new
sequence has at most \(c\ell\) members and ends in \(E\).

Each term occurring in a copied member occurs in \(D\).  Each term \(b\)
occurring in \(\theta_iT_\theta\) has, by \(\operatorname{lh}\)-directness,
\(\operatorname{lh}(b)\leq\operatorname{lh}(c')\) for some term \(c'\)
occurring in \(\theta E_\rho\) or \(\theta S_\rho\), which are members of
\(D\).
\end{proof}
\end{annexitem}
